\documentclass[a4paper,11pt]{article}

% =============================================================================
%  Charte graphique des documents LaTeX du projet (gabarit commenté).
%  Le préambule fixe la mise en forme ; le corps montre chaque élément une fois
%  et se remplace par le contenu du document. Un document reprend le préambule
%  sans modification, hors des passages bornés par les deux lignes
%  « % >>> À ADAPTER : <objet> » et « % <<< À ADAPTER » (comparaison par
%  script : à créer, ADR 0003 annoté le 7 novembre 2026). Le bloc « Macros du dépôt »,
%  en fin de préambule, porte \src et les rubriques de la fiche modèle.
%  Version sans : babel-french, caption, csquotes, dsfont, enumitem, float,
%  mathrsfs, mdframed, newtxtext, pdflscape, tikz/pgf, titlesec, tocloft.
%  Moteur : pdfLaTeX. Compatible avec les distributions anciennes
%  (MiKTeX 2.9, TeX Live 2012 et suivantes) : aucun crochet du noyau récent.
% =============================================================================

\usepackage[utf8]{inputenc}
\usepackage[T1]{fontenc}
% \headheight est laissé à 12 pt : fancyhdr le signale et l'agrandit lui-même
% (avertissement « \headheight is too small »), exactement comme dans le
% document source ; le fixer ici décalerait l'en-tête ou le texte de 1,6 pt.
\usepackage[left=2.5cm, right=2.5cm, top=2.5cm, bottom=2.5cm]{geometry}
\usepackage{setspace}
\usepackage{graphicx}
\usepackage{etoolbox}
\usepackage{keyval}
\usepackage{fancyhdr}
\usepackage{amsmath,amssymb,amsfonts}
\usepackage{url}
\usepackage[table]{xcolor}  % [table] : \rowcolors sur les anciennes versions
\usepackage{array}
\usepackage{booktabs}
\usepackage{colortbl}
\usepackage{longtable}
\usepackage{lscape}   % remplace pdflscape (rotation à l'écran : voir plus bas)
\usepackage{framed}   % remplace mdframed (encadré sécable)
\usepackage{atbegshi}  % crochets de sortie de page (remplace \AddToHook)
\usepackage[unicode]{hyperref}

% ------------------------------------------------------------- polices -------
% newtxtext composait le texte en TeX Gyre Termes X et le code en txtt ; les
% mathématiques restaient en Computer Modern. On retrouve la même chose sans
% newtxtext : TeX Gyre Termes (tgtermes) pour le texte, repli sur le Times de
% PSNFSS (ptm) s'il est absent ; txtt (txfonts) pour \ttfamily.
\IfFileExists{tgtermes.sty}
  {\usepackage{tgtermes}}
  {\renewcommand{\rmdefault}{ptm}}
\IfFileExists{t1txtt.fd}{\renewcommand{\ttdefault}{txtt}}{}

% ---------------------------------------------------------------- charte ----
% Bleu financier pour la structure, vert pour l'accent et les sorties, deux
% gris de tableau.
\definecolor{maincolor}{RGB}{0,70,140}
\definecolor{greenline}{RGB}{0,180,80}
\definecolor{tableheader}{RGB}{200,220,250}
\definecolor{tablerow}{RGB}{245,245,245}
\definecolor{financeblue}{RGB}{0,70,140}

% >>> À ADAPTER : métadonnées du document (page de garde, pied de page, PDF)
% Dans un guide, \datedoc est une date figée (jamais \today) et \statutdoc
% vaut Projet ou Final ; \datedroit est la date de référence du droit.
\newcommand{\annee}{2026}
\newcommand{\datedoc}{9 novembre 2026}
\newcommand{\versiondoc}{0.1}
\newcommand{\statutdoc}{Projet}
\newcommand{\titredoc}{Documentation du package provie}
\newcommand{\auteurdoc}{Projet outil\_prov\_vie}
\newcommand{\datedroit}{Sans objet}

\hypersetup{colorlinks=true, linkcolor=financeblue, urlcolor=financeblue,
            citecolor=financeblue, pdftitle={Documentation du package provie},
            pdfauthor={\auteurdoc}}
% <<< À ADAPTER

\makeatletter

% ================================================ français sans babel ========
% Reprise, à la main, des effets de babel-french visibles dans le document
% source.
% (1) Césure française (motifs chargés dans le format pdfLaTeX).
\ifdefined\l@french
  \language=\l@french
\else
  \PackageWarningNoLine{charte}{Motifs de césure français absents :
    césure anglaise utilisée}
\fi
\lefthyphenmin=2 \righthyphenmin=2
% (2) Espacement « à la française » après les points.
\frenchspacing
% (3) Libellés.
\renewcommand{\contentsname}{Table des matières}
\renewcommand{\listfigurename}{Table des figures}
\renewcommand{\listtablename}{Liste des tableaux}
\renewcommand{\refname}{Références}
\renewcommand{\indexname}{Index}
\renewcommand{\figurename}{Figure}
\renewcommand{\tablename}{Table}
\renewcommand{\partname}{Partie}
\renewcommand{\appendixname}{Annexe}
\renewcommand{\abstractname}{Résumé}
% (4) Date en français (« 1er janvier 2026 »).
\renewcommand{\today}{\number\day\ifnum\day=1 \textsuperscript{er}\fi\space
  \ifcase\month\or janvier\or février\or mars\or avril\or mai\or juin\or
  juillet\or août\or septembre\or octobre\or novembre\or décembre\fi
  \space\number\year}
% (5) Premier paragraphe indenté après un titre.
\let\@afterindentfalse\@afterindenttrue
\@afterindenttrue
% (6) Listes : espacement vertical resserré de babel-french, appliqué à tous
%     les environnements construits sur \list.
\let\charte@listORI\list
\newdimen\charte@parskip
\def\charte@listV{%
  \setlength{\topsep}{0.8ex plus 0.4ex minus 0.4ex}%
  \setlength{\partopsep}{0.4ex plus 0.2ex minus 0.2ex}%
  \charte@parskip=\parskip
  \addtolength{\topsep}{-\charte@parskip}%
  \addtolength{\partopsep}{\charte@parskip}%
  \setlength{\itemsep}{0.4ex plus 0.2ex minus 0.2ex}%
  \setlength{\parsep}{0.4ex plus 0.2ex minus 0.2ex}%
  \ifdim\charte@parskip>\z@
    \setlength{\parsep}{\charte@parskip}%
    \addtolength{\itemsep}{-\charte@parskip}%
  \fi}
\def\list#1#2{\charte@listORI{#1}{\charte@listV #2}}
% (7) Ponctuation haute : sans babel-french, l'espace qui précède « : » et
%     l'espace fine qui précède « ; ! ? » ne sont plus ajoutées. On saisit
%     « ~: » (espace insécable) et « \fsp; », « \fsp! », « \fsp? » (espace
%     fine insécable, de même largeur que celle de babel-french).
\newcommand{\fsp}{\unskip\nobreak\hskip .5\fontdimen2\font\relax}
% (8) Puces en tiret cadratin à tous les niveaux.
\renewcommand{\labelitemi}{\textemdash}
\renewcommand{\labelitemii}{\textemdash}
\renewcommand{\labelitemiii}{\textemdash}
\renewcommand{\labelitemiv}{\textemdash}

% ================================================ listes sans enumitem =======
% itemize et enumerate acceptent un argument optionnel à clés, dont la
% syntaxe est celle d'enumitem pour les clés usuelles :
%   \begin{itemize}[itemsep=3pt] ... \end{itemize}
% Clés reconnues : itemsep, topsep, parsep, partopsep, leftmargin,
% labelsep, listparindent.
\define@key{charteliste}{itemsep}{\setlength{\itemsep}{#1}}
\define@key{charteliste}{topsep}{\setlength{\topsep}{#1}}
\define@key{charteliste}{parsep}{\setlength{\parsep}{#1}}
\define@key{charteliste}{partopsep}{\setlength{\partopsep}{#1}}
\define@key{charteliste}{leftmargin}{\setlength{\leftmargin}{#1}}
\define@key{charteliste}{labelsep}{\setlength{\labelsep}{#1}}
\define@key{charteliste}{listparindent}{\setlength{\listparindent}{#1}}

\renewenvironment{itemize}[1][]{%
  \ifnum\@itemdepth>\thr@@\@toodeep\else
    \advance\@itemdepth\@ne
    \edef\@itemitem{labelitem\romannumeral\the\@itemdepth}%
    \expandafter\list\csname\@itemitem\endcsname
      {\def\makelabel##1{\hss\llap{##1}}\setkeys{charteliste}{#1}}%
  \fi}{\endlist}

\renewenvironment{enumerate}[1][]{%
  \ifnum\@enumdepth>\thr@@\@toodeep\else
    \advance\@enumdepth\@ne
    \edef\@enumctr{enum\romannumeral\the\@enumdepth}%
    \expandafter\list\csname label\@enumctr\endcsname
      {\usecounter\@enumctr\def\makelabel##1{\hss\llap{##1}}%
       \setkeys{charteliste}{#1}}%
  \fi}{\endlist}

% ================================================ légendes sans caption ======
% Étiquette en gras, couleur maincolor, petites capitales ; séparateur
% « : » précédé d'une espace insécable, comme dans le document source.
\renewcommand{\fnum@figure}{{\scshape\figurename}\nobreakspace\thefigure}
\renewcommand{\fnum@table}{{\scshape\tablename}\nobreakspace\thetable}
\newcommand{\charte@etiquette}[1]{%
  {\bfseries\color{maincolor}#1\nobreakspace:}\space}
% Comme avec caption : légende composée en interligne simple, entre deux
% étais, sans ressort interligne au-dessus ; centrée si elle tient sur une
% ligne, justifiée sinon.
\long\def\@makecaption#1#2{%
  \vskip\abovecaptionskip
  \sbox\@tempboxa{\normalfont\charte@etiquette{#1}#2}%
  \nointerlineskip
  \ifdim\wd\@tempboxa>\hsize
    \hb@xt@\hsize{\vbox{\setstretch{1}\normalfont
      \parindent\z@ \leftskip\z@ \rightskip\z@ \parfillskip\z@\@plus1fil
      \strut\charte@etiquette{#1}#2\strut\par}}%
  \else
    \global\@minipagefalse
    \hb@xt@\hsize{\hfil\setstretch{1}\normalfont\strut\box\@tempboxa\hfil}%
  \fi
  \vskip\belowcaptionskip}

% Figures et tableaux « ici » : reprise exacte de la position [H] du
% package float (boîte non flottante, \intextsep avant et après).
\newsavebox\charte@boiteH
\def\charte@debutH#1{%
  \def\@captype{#1}%
  \global\setbox\charte@boiteH\color@vbox\normalcolor
    \vbox\bgroup\hsize\columnwidth\@parboxrestore
      \@floatboxreset\@setnobreak\ignorespaces}
\def\charte@finH{%
  \@endfloatbox
  \vskip\intextsep\box\charte@boiteH\vskip\intextsep\relax}
\newenvironment{figureH}{\charte@debutH{figure}}{\charte@finH}
\newenvironment{tableH}{\charte@debutH{table}}{\charte@finH}

% ================================================ en-têtes et pieds ==========
\renewcommand{\baselinestretch}{1.08}
\pagestyle{fancy}
\fancyhf{}
\fancyhead[L]{\color{maincolor}\nouppercase{\leftmark}}
\renewcommand{\headrulewidth}{0pt}
\renewcommand{\footrulewidth}{0.4pt}
\fancyfoot[R]{\color{greenline}\textbf{page \thepage}}
\fancyfoot[L]{\color{maincolor}\titredoc}
\fancypagestyle{plain}{\fancyhf{}%
  \renewcommand{\headrulewidth}{0pt}\renewcommand{\footrulewidth}{0.4pt}%
  \fancyfoot[R]{\color{greenline}\textbf{page \thepage}}%
  \fancyfoot[L]{\color{maincolor}\titredoc}}

% ================================================ titres sans titlesec =======
% Bandeau : filet bleu pleine largeur, titre vert en dessous. Utilisé seul
% pour les pages non numérotées (objet, sommaire) et par \section.
\newcommand{\titrebandeau}[1]{%
  \noindent{\color{financeblue}\rule{\linewidth}{2.5pt}}
  \par\vspace{5mm}
  {\color{greenline}\raggedright\fontsize{19}{23}\bfseries\selectfont #1\par}
  \vspace{6mm}}

% Numéros de titre : aucun pour \section (le numéro ne figure qu'au
% sommaire et dans l'en-tête), 0,7 em puis 0,6 em d'écart ensuite.
\def\@seccntformat#1{\csname charte@num@#1\endcsname}
\def\charte@num@section{}
\def\charte@num@subsection{\strut\thesubsection\hspace{0.7em}}
\def\charte@num@subsubsection{\strut\thesubsubsection\hspace{0.6em}}
\def\charte@num@paragraph{\theparagraph\quad}
\def\charte@num@subparagraph{\thesubparagraph\quad}

% \@startsection passe le titre, entre accolades, à la dernière commande de
% son argument de style. Pour \section, cette commande est \charte@bandeau :
% une ligne vide (étai) au-dessus du bandeau et une autre en dessous, et un
% étai dans chaque ligne, comme la forme « display » de titlesec avec
% étiquette vide.
\newcommand{\charte@bandeau}[1]{%
  \parindent\z@ \everypar{\strut}%
  {\strut\par}\nobreak
  \titrebandeau{#1}\nobreak
  \kern\z@\nobreak\strut\par}
% Titres en ligne : un étai en fin de titre (le premier est dans le numéro).
\newcommand{\charte@etai}[1]{\charte@etai@#1}
\def\charte@etai@#1\@@par{#1\kern\z@\strut\@@par}
% Réserve avant un titre : si le bas de la page ne peut plus recevoir le titre
% et trois lignes de texte, le titre passe à la page suivante (M41 : aucun
% titre en bas de page avec moins de trois lignes sous lui). La réserve
% compte le titre lui-même et ses espaces (mesurés : bandeau de section
% 7,4 lignes, sous-section 2,9 lignes), trois lignes et le cadre d'un
% encadré qui suivrait : 12 lignes pour \section, 7 pour les autres. Logique de la
% macro \Needspace* du package needspace, sans le charger : la place restante
% est mesurée (\pagegoal - \pagetotal, après une pénalité nulle qui met la
% page à jour) et le saut n'est forcé que si elle manque. La forme élastique
% de \needspace, employée jusqu'ici, créait un point de coupure de faible
% coût que TeX préférait à toute autre coupure dès qu'un bloc insécable
% suivait (tableau, encadré), même avec la page aux trois quarts vide. Sans
% effet juste après un autre titre (\if@nobreak), où elle créerait un point
% de coupure sous ce titre, ni dans une boîte (\ifinner : minipage, \parbox),
% où la page en cours ne la concerne pas.
\newcommand{\charte@reserve}[1]{\par\ifinner\else\if@nobreak\else
  \penalty\z@ % met à jour \pagetotal
  \begingroup\setlength{\dimen@}{#1}%
  \dimen@ii\pagegoal \advance\dimen@ii-\pagetotal
  \ifdim\pagegoal=\maxdimen\else\ifdim\dimen@>\dimen@ii
    \vfil\penalty-\@M % place insuffisante : page suivante
  \fi\fi\endgroup\fi\fi}
\renewcommand{\section}{\charte@reserve{12\baselineskip}%
  \@startsection{section}{1}{\z@}%
  {\z@}{1sp}{\normalfont\charte@bandeau}}
% Espacements : ceux de titlesec, soit 3,25 ex (+1 ex, -0,2 ex) avant et
% 1,5 ex (+0,2 ex) après, l'ex étant celui que titlesec mémorise au
% chargement (Computer Modern à 11 pt, soit 4,71341 pt).
\newdimen\charte@ex \charte@ex=4.71341pt
\newskip\charte@avanttitre
\charte@avanttitre=-3.25\charte@ex \@plus -1\charte@ex \@minus -.2\charte@ex
\newskip\charte@aprestitre
\charte@aprestitre=1.5\charte@ex \@plus .2\charte@ex
\renewcommand{\subsection}{\charte@reserve{7\baselineskip}%
  \@startsection{subsection}{2}{\z@}%
  {\charte@avanttitre}{\charte@aprestitre}%
  {\normalfont\large\bfseries\color{financeblue}\charte@etai}}
\renewcommand{\subsubsection}{\charte@reserve{7\baselineskip}%
  \@startsection{subsubsection}{3}{\z@}%
  {\charte@avanttitre}{\charte@aprestitre}%
  {\normalfont\normalsize\bfseries\color{financeblue}\charte@etai}}

% ================================================ sommaire sans tocloft ======
% Sections : gras bleu, numéro sur 2,2 em, numéro de page bleu (maigre),
% sans points de conduite. Sous-sections : retrait 1,5 em, numéro sur 2,3 em,
% points de conduite. L'em est celui que tocloft mémorisait au chargement
% (10,88788 pt), pour un sommaire identique au document source.
\newdimen\charte@em \charte@em=10.88788pt
\renewcommand*{\l@section}[2]{%
  \ifnum\c@tocdepth>\z@
    \addpenalty\@secpenalty
    \addvspace{1.0\charte@em \@plus\p@}%
    \setlength\@tempdima{2.2\charte@em}%
    \begingroup
      \parindent\z@ \rightskip\@pnumwidth
      \parfillskip -\@pnumwidth
      \leavevmode
      \advance\leftskip\@tempdima
      \hskip -\leftskip
      {\bfseries\color{financeblue}#1}\nobreak\hfil\nobreak
      \hb@xt@\@pnumwidth{\hss\color{financeblue}#2}\par
    \endgroup
  \fi}
\renewcommand*{\l@subsection}{\@dottedtocline{2}{1.5\charte@em}{2.3\charte@em}}
\setcounter{tocdepth}{2}

% ================================================ environnements de texte ====
\newcommand{\code}[1]{\texttt{\small #1}}
\newcommand{\refl}[1]{{\small\textcolor{maincolor!70!black}{\textit{Référence.} #1}}\par\smallskip}

% Encadré de citation réglementaire : filet vertical bleu en marge.
\newcommand{\reglement}[2]{%
  \par\vspace{2mm}\noindent\begin{minipage}{\linewidth}
  \small\setlength{\parindent}{0pt}
  \textcolor{financeblue}{\vrule width 2pt}\hspace{4mm}%
  \begin{minipage}{\dimexpr\linewidth-7mm\relax}
  \textcolor{financeblue}{\textbf{#1}}
  \par\vspace{2pt}
  {\itshape #2}
  \end{minipage}
  \end{minipage}\par\vspace{3mm}}

% Encadré générique : fond gris, marge intérieure 6 mm à gauche et 10 pt à
% droite, 4 pt en haut et en bas, 3 mm avant (sauf après un titre) et
% après ; premier paragraphe non
% indenté. Sécable entre pages (framed), selon \charte@coupe ci-dessous. Le
% cadre garde la profondeur de la boîte d'essai de framed : sa hauteur
% mesurée est de 8 pt (elle valait -13,68 pt quand la profondeur était
% annulée, et framed taillait le premier morceau trop grand pour la page).
\newenvironment{encadre}{%
  \def\FrameCommand##1{{\setlength{\fboxsep}{\z@}%
    \colorbox{tablerow}{\hspace*{6mm}%
      \vbox{\vskip4pt\hbox{##1}\vskip4pt}\hspace*{10pt}}}}%
  \@ifundefined{OuterFrameSep}{}{\setlength{\OuterFrameSep}{\z@}}%
  \topsep\z@ \partopsep\z@ % framed antérieur à 0.96 : espacement via \topsep
  \par
  \if@nobreak\charte@sanscoupure\else\vskip 3mm\fi
  \nointerlineskip
  \MakeFramed{\let\penalty\charte@penalite@ori
    \advance\hsize-\width \FrameRestore
    \everypar{{\setbox\z@\lastbox}\everypar{}}}}%
  {\global\let\fb@put@frame\charte@coupe % hors du groupe de la boîte
   \global\charte@premiertrue
   \endMakeFramed
   \global\let\fb@put@frame\charte@putframe
   \par
   \vskip 3mm\prevdepth 4pt
   \global\@nobreakfalse\global\@endpetrue}

% Coupe de l'encadré (E3.7, M41), décidée à chaque placement d'un morceau,
% avant la coupe de framed. Premier placement : un encadré court (hauteur
% totale au plus \charte@encadrecourt) qui ne tient pas dans la place
% restante est posé entier, et TeX coupe la page au dernier point de coupure
% qui le précède : juste avant lui, ou avant le titre qu'il suit (aucun point
% de coupure entre un titre et l'encadré : \charte@sanscoupure). Encadré
% long : coupé, chaque morceau gardant au moins trois lignes, à chaque page
% (page suivante si la place restante en tient moins ; page écourtée si le
% reste en tiendrait moins). Exception : si des flottants attendent la page
% suivante (\@deferlist) et que l'encadré court ne tient pas avec eux sur une
% page, posé entier il déborderait ; il est alors coupé comme un encadré
% long (M200). De même s'il commence en tête de page : il ne peut aller plus
% loin, et la page peut déjà porter des flottants (M202). Un encadré coupé trop
% court pour donner deux morceaux de trois lignes (hauteur du contenu, cadre
% déduit, sous 5,5 interlignes : six lignes ou moins ; sept lignes se coupent
% en trois et quatre) passe entier à la page suivante. Limites : tous les
% flottants en attente sont comptés, même ceux qui n'iront pas sur la page
% suivante ([p], [b] trop grand, \topnumber dépassé) : coupe parfois
% inutile, avertie par le contrôle 13 ; sur une page qui porte un flottant
% en tête, la coupe de framed peut déborder d'environ 1 pt (Overfull \vbox
% mesuré jusqu'à 1,03 pt, toléré en AVERT par le contrôle 3 jusqu'à 1,5 pt,
% M202).
\newdimen\charte@encadrecourt \charte@encadrecourt=.5\textheight
\newif\ifcharte@premier
\let\charte@putframe\fb@put@frame
\let\charte@penalite@ori\penalty
% Pénalités 9999 de \MakeFramed (mise à jour de \pagetotal) rendues
% insécables juste après un titre ; \penalty est rétabli dans la boîte.
\def\charte@sanscoupure{\def\penalty{\afterassignment\charte@penalite\count@}}
\def\charte@penalite{\ifnum\count@=9999 \count@\@M\fi\charte@penalite@ori\count@}
\def\charte@coupe#1#2{%
  \let\charte@suite\charte@putframe
  \ifinner\else\ifdim\pagegoal=\maxdimen\else
    \fb@sizeofframe#1%
    \dimen@\pagegoal \advance\dimen@-\pagetotal \advance\dimen@-\fb@frh
    \advance\dimen@-\pageshrink \advance\dimen@-\pagedepth
    \dimen@ii\ht\@tempboxa
    \ifdim\dimen@ii>\dimen@
      \advance\dimen@ii\fb@frh
      \ifcharte@premier\ifdim\dimen@ii>\charte@encadrecourt\else
        \charte@flottants % hauteur des flottants en attente (M200)
        \advance\dimen@ii\@tempdima
        \@tempdimb\topskip \advance\@tempdimb\p@
        \ifdim\pagetotal>\@tempdimb % la page porte déjà du texte
          \ifdim\dimen@ii>\textheight\else % tient avec eux : posé entier
            \let\charte@suite\charte@entier
          \fi
        \fi % en tête de page (flottants déjà placés) : coupé (M202)
        \advance\dimen@ii-\@tempdima
      \fi\fi
      \advance\dimen@ii-\fb@frh
      \ifx\charte@suite\charte@entier\else
        \ifdim\dimen@<3\baselineskip
          \eject\fb@adjheight
        \else
          \@tempdimb\dimen@ii \advance\@tempdimb-\fb@frh
          \advance\dimen@ii-\dimen@
          \ifdim\dimen@ii<3\baselineskip
            \ifdim\@tempdimb<5.5\baselineskip
              % trop court pour deux morceaux de trois lignes : page suivante
              \eject\fb@adjheight
            \else
              \advance\dimen@ii-3\baselineskip
              \advance\dimen@ii-\fb@frh % marge : rembourrage du cadre
              \advance\dimen@ii\pagegoal \pagegoal\dimen@ii % toujours global
            \fi
          \fi
        \fi
      \fi
    \fi
  \fi\fi
  \global\charte@premierfalse
  \charte@suite#1#2}
% Hauteur (\@tempdima) des flottants qui attendent la page suivante
% (\@deferlist), avec leurs espaces ; zéro s'il n'y en a pas.
\def\charte@flottants{\@tempdima\z@
  \ifx\@deferlist\@empty\else
    \begingroup
      \def\@elt##1{\global\advance\@tempdima\ht##1%
        \global\advance\@tempdima\dp##1\global\advance\@tempdima\floatsep}%
      \@deferlist
    \endgroup
    \advance\@tempdima\textfloatsep \advance\@tempdima-\floatsep
  \fi}
% Encadré court posé entier : la page déborde et TeX la coupe avant lui.
\def\charte@entier#1#2{\fb@putboxa#1\fb@afterframe}

% Fiches : un environnement, deux jeux de rubriques fixes. Fiche méthode
% (corps du guide) : six rubriques, ouvertes par \Principe, ci-dessous ;
% fiche modèle (annexe technique) : cinq rubriques, ouvertes par \Cadre, au
% bloc « Macros du dépôt ».
\newcounter{ficheid}
\newenvironment{fiche}[1]{%
  \stepcounter{ficheid}%
  \phantomsection
  \subsection{#1}%
  \label{fiche:\theficheid}%
}{%
  \end{rubriques}\vspace{0.5em}%
}
% Liste des rubriques : aucun retrait ; l'intitulé fait partie du
% paragraphe (ses espaces se justifient avec la ligne, comme le style
% « unboxed » d'enumitem), suivi d'un écart fixe \labelsep.
\newenvironment{rubriques}{%
  \list{}{\setlength{\leftmargin}{\z@}\setlength{\labelwidth}{\z@}%
          \setlength{\itemindent}{\z@}%
          \setlength{\itemsep}{3pt}\setlength{\topsep}{3pt}%
          \setlength{\parsep}{2pt}%
          \let\makelabel\@gobble}%
  \clubpenalties 3 10000 10000 0
  \@beginparpenalty\@M \global\charte@rubriques\z@}{\endlist}
% Un titre de rubrique n'est jamais seul en bas de page : il fait partie de
% son premier paragraphe, dont les deux premières coupures sont interdites
% (\clubpenalties, liste des rubriques), et une liste qui le suit ne s'en
% sépare pas (\@beginparpenalty) ; aucune réserve de place, qui laisserait
% orphelin le titre de la fiche ; aucune coupure avant la deuxième ni la
% troisième rubrique, pour que le titre de la fiche garde trois lignes sous
% lui même quand les premières rubriques tiennent sur une ligne. Le réglage
% \clubpenalties vaut pour tous les paragraphes de la fiche, listes
% imbriquées comprises : aucun ne laisse moins de trois lignes en bas de page
% (un paragraphe court passe alors entier à la page suivante).
% Limite : un titre de rubrique suivi d'une liste peut garder une seule
% puce sous lui en bas de page (jamais aucune) ; la règle des trois lignes
% (M41) est celle des titres de section (M202).
\newcount\charte@rubriques
\newcommand{\rubrique}[1]{%
  \global\advance\charte@rubriques\@ne
  \ifnum\charte@rubriques<4 \@itempenalty\@M\else\@itempenalty-\@lowpenalty\fi
  \item[]\@itempenalty-\@lowpenalty % listes de la rubrique : sécables
  \textcolor{financeblue}{\textbf{#1}}\hskip\labelsep\ignorespaces}
\newcommand{\Principe}{%
  \begin{rubriques}%
  \rubrique{1. Principe et modèle de référence.}}
\newcommand{\Donnees}{\rubrique{2. Données requises.}}
\newcommand{\Conditions}{\rubrique{3. Conditions d'usage.}}
\newcommand{\Forces}{\rubrique{4. Forces.}}
\newcommand{\Limites}{\rubrique{5. Limites et pièges.}}
\newcommand{\Referentiels}{\rubrique{6. Référentiels concernés.}}

% Titre de partie : un filet bleu, « Partie » en petites capitales, titre
% vert.
\newcommand{\partie}[2]{%
  \clearpage
  \vspace*{1cm}
  \noindent{\color{financeblue}\rule{\linewidth}{2pt}}\par\vspace{3mm}
  {\color{financeblue}\raggedright\normalsize\bfseries\scshape #1\par}
  \vspace{1.5mm}
  {\color{greenline}\raggedright\fontsize{20}{24}\bfseries\selectfont #2\par}
  \vspace{5mm}
  \addcontentsline{toc}{section}{\protect\numberline{}\textsc{#1}~: #2}}

% ================================================ page de garde sans TikZ ====
% \charte@surpage{<matériel vertical>} superpose, sur la prochaine page
% expédiée, un contenu placé depuis le coin supérieur gauche du papier
% (atbegshi : fonctionne avec les noyaux LaTeX anciens comme récents).
\newcommand{\charte@surpage}[1]{%
  \AtBeginShipoutNext{%
    \global\setbox\AtBeginShipoutBox\vbox{\offinterlineskip
      \vbox to\z@{\kern-1in\kern-\voffset
        \hbox to\z@{\kern-1in\kern-\hoffset
          \vbox to\z@{#1\vss}\hss}\vss}%
      \box\AtBeginShipoutBox}}}
% Trois filets : bleu de 3,2 à 3,28 cm du haut avec amorce verte de 5 cm,
% bleu de 3,6 à 3,53 cm du bas, entre les marges de 2,5 cm.
% \leavevmode ouvre le paragraphe comme le faisait la figure TikZ du document
% source : la mention d'en-tête se place ainsi au-dessus du filet, en retrait.
\newcommand{\filetsgarde}{%
  \charte@surpage{%
    \kern3.2cm
    \moveright2.5cm\hbox{\color{financeblue}%
      \rule{\dimexpr\paperwidth-5cm\relax}{0.08cm}}%
    \kern-0.08cm
    \moveright2.5cm\hbox{\color{greenline}\rule{5cm}{0.08cm}}%
    \kern\dimexpr\paperheight-3.6cm-3.28cm\relax
    \moveright2.5cm\hbox{\color{financeblue}%
      \rule{\dimexpr\paperwidth-5cm\relax}{0.07cm}}}%
  \leavevmode}

% ================================================ pages en paysage ===========
% lscape tourne le contenu ; comme pdflscape, on ajoute /Rotate 90 aux pages
% concernées pour qu'elles s'affichent à l'endroit à l'écran.
\newif\ifcharte@paysage
\apptocmd{\landscape}{\global\charte@paysagetrue}{}{}
\apptocmd{\endlandscape}{\global\charte@paysagefalse}{}{}
\AtBeginShipout{%
  \ifcharte@paysage\global\pdfpageattr{/Rotate 90}%
  \else\global\pdfpageattr{}\fi}

% ================================================ Macros du dépôt ============
% Propres au dépôt de guides, identiques dans tous les documents.
%
% (1) Statuts de vérification des sources. statuts_sources.tex est généré
%     depuis le registre des sources (sources/registre.csv) par
%     outils/controles.sh, et trouvé par TEXINPUTS (répertoire sources/) ;
%     il contient une ligne \charte@statut{ID}{statut} par identifiant, le
%     statut valant verifie, incertain, nontrouve ou absent.
\newif\ifcharte@statuts
\newcommand{\charte@statut}[2]{\@namedef{charte@src@#1}{#2}}
\def\charte@verifie{verifie}
\IfFileExists{statuts_sources.tex}
  {\input{statuts_sources.tex}\charte@statutstrue}
  {\PackageWarningNoLine{charte}{statuts_sources.tex introuvable :
     statuts de vérification des sources non calculés}}
% (2) \src{ID} : identifiant du registre des sources, placé dans le premier
%     argument de \reglement ou dans l'argument de \refl, en fin de texte.
%     Statut vérifié : rien d'imprimé ; incertain, non trouvé ou absent : une
%     marque † en exposant, accolée au mot qui précède ; identifiant inconnu :
%     [ID ?] en gras et un avertissement « Source inconnue au registre ».
\DeclareRobustCommand{\src}[1]{%
  \ifcharte@statuts
    \@ifundefined{charte@src@#1}%
      {\textbf{[#1~?]}%
       \PackageWarning{charte}{Source inconnue au registre : #1}}%
      {\unskip
       \expandafter\ifx\csname charte@src@#1\endcsname\charte@verifie
       \else\textsuperscript{\dag}\fi}%
  \fi
  \protected@write\@auxout{}{\string\charte@srcpage{#1}{\thepage}}}
% (2 bis) Folios des \src, pour l'annexe des sources citées.
\newcommand{\charte@srcpage}[2]{%
  \@ifundefined{charte@pg@#1}%
    {\global\@namedef{charte@pg@#1}{\hyperpage{#2}}%
     \global\@namedef{charte@dp@#1}{#2}}%
    {\def\charte@tmp{#2}%
     \expandafter\ifx\csname charte@dp@#1\endcsname\charte@tmp\else
       \expandafter\g@addto@macro\csname charte@pg@#1\endcsname{, \hyperpage{#2}}%
       \global\@namedef{charte@dp@#1}{#2}\fi}}
\AtEndDocument{\let\charte@srcpage\@gobbletwo}
\newcommand{\srcfolios}[1]{%
  \@ifundefined{charte@pg@#1}%
    {\textbf{??}\G@refundefinedtrue
     \PackageWarning{charte}{Folios de la source #1 non definis}}%
    {\@nameuse{charte@pg@#1}}}
% (3) Rubriques de la fiche modèle (annexe technique), sur l'environnement
%     fiche : cinq rubriques fixes, ouvertes par \Cadre.
\newcommand{\Cadre}{%
  \begin{rubriques}%
  \rubrique{1. Cadre et notations.}}
\newcommand{\Hypotheses}{\rubrique{2. Hypothèses.}}
\newcommand{\Utilisation}{\rubrique{3. Conditions d'utilisation.}}
\newcommand{\Formules}{\rubrique{4. Formules et statut.}}
\newcommand{\LimitesModele}{\rubrique{5. Limites.}}
% (4) Bibliographie (BibTeX, style plain, base references, soit
%     sources/references.bib trouvé par BIBINPUTS) : \bibliographiedepot, en
%     fin de document, ne la compose que si le document cite une entrée par
%     \cite ou \nocite ; sans citation, bibtex s'arrêterait en erreur
%     (« I found no \citation commands »).
\newif\ifcharte@citations
\pretocmd{\@citex}{\global\charte@citationstrue}{}
  {\PackageWarningNoLine{charte}{\string\@citex\space non modifiée :
     bibliographie non composée}}
\pretocmd{\nocite}{\global\charte@citationstrue}{}
  {\PackageWarningNoLine{charte}{\string\nocite\space non modifiée :
     bibliographie non composée}}
\newcommand{\bibliographiedepot}{%
  \ifcharte@citations
    \bibliographystyle{plain}%
    \bibliography{references}%
  \fi}
% (5) Annexe des sources citées. \sourcescitees charge sources_citees.tex,
%     tableau long généré à chaque compilation par
%     outils/controles/sources_citees.R depuis le registre des sources et les
%     \src du document, et trouvé par TEXINPUTS (répertoire de construction) ;
%     l'annexe ne se saisit jamais à la main.
\newcommand{\sourcescitees}{%
  \IfFileExists{sources_citees.tex}
    {\input{sources_citees.tex}}
    {\PackageWarningNoLine{charte}{sources_citees.tex introuvable :
       annexe des sources citées non générée}}}

\makeatother

\emergencystretch=1.5em

\begin{document}

% ============================================================ page de garde ==
\thispagestyle{empty}
\filetsgarde
\vspace*{2.1cm}
\noindent
{\color{financeblue}\normalsize\bfseries\scshape Édition \annee}

\vspace{1.7cm}
\noindent
{\color{financeblue}\fontsize{28}{34}\bfseries\selectfont
 Documentation\\ du package provie\par}

\vspace{1.1cm}
\noindent
{\color{black}\Large
 Recalcul des provisions d'assurance vie et de prévoyance\\
 Modules, fonctions et concordance avec le code\par}

\vspace{0.7cm}
\noindent
{\color{black}\normalsize\itshape
 Garanties vie et prévoyance assurables en France, organismes d'assurance\\
 Comptes sociaux et Solvabilité~II\par}

\vfill
\noindent
{\footnotesize\color{black}
 Établie le \datedoc\quad\textbar\quad Version \versiondoc\quad\textbar\quad
 Statut~: \statutdoc\quad\textbar\quad Auteur~: \auteurdoc\par}
\vspace{2.6cm}
\clearpage

% ================================================================= objet ====
\titrebandeau{Objet du document}
\markboth{Objet du document}{}

\noindent
Ce document est la \textbf{documentation de fond} du package R
\code{provie}, moteur de l'outil \code{outil\_prov\_vie}. Il décrit
l'architecture du package, puis, module par module, son rôle et ses
fonctions. Il s'adresse à un relecteur externe, qui doit pouvoir comprendre
ce que fait le package sans lire le code.

\medskip\noindent
Le document suit l'état du code. Chaque module a son tableau de fonctions,
vide tant que le module n'en a pas\fsp; ces tableaux sont comparés au code
par une batterie de tests (section~\ref{sec:concordance}). La documentation ne porte aucun
chiffre issu d'une exécution~: les résultats d'un calcul sont restitués par
le rapport figé et la fiche de calcul.

\vspace{6mm}
\noindent
\begin{center}
\small
\rowcolors{2}{tablerow}{white}
\begin{tabular}{@{}p{5.2cm}p{9.6cm}@{}}
\rowcolor{tableheader}
\textcolor{financeblue}{\textbf{Élément}} & \textcolor{financeblue}{\textbf{Référence}} \\
Cahier des charges & \code{docs/exigences.md} \\
Vocabulaire & \code{CONTEXT.md} \\
Décisions d'architecture & \code{docs/adr/} \\
Conventions de rédaction & \code{docs/latex/CONVENTIONS.md} \\
Code du package & \code{R/}, un fichier par module \\
Batterie de concordance & \code{tests/concordance\_doc\_moteur.R} \\
\end{tabular}
\end{center}

\clearpage
\titrebandeau{Sommaire}
\markboth{Sommaire}{}
% \tableofcontents composerait son propre titre, en doublon avec le bandeau :
% on n'appelle que \@starttoc, qui compose les entrées seules.
\makeatletter
\@starttoc{toc}
\makeatother
\newpage


\partie{Partie I}{Présentation}

% =============================================================================
\section{Périmètre}
\label{sec:perimetre}
% =============================================================================

Le package \code{provie} recalcule, à une date d'arrêté déclarée, les
provisions des garanties d'assurance vie et de prévoyance assurables en
France par un organisme d'assurance (société d'assurance, mutuelle,
institution de prévoyance). Il les calcule sous deux référentiels, les
comptes sociaux et la meilleure estimation de Solvabilité~II, et restitue le
tableau de passage de l'un à l'autre.

\begin{encadre}
\textbf{Portée.} Le package est un outil de \textbf{validation actuarielle}
(second regard de la fonction actuarielle, revue externe), jamais un calcul
de production. Les provisions inscrites par l'organisme d'assurance ne sont
pas une entrée, et le package ne rend aucun verdict de suffisance. IFRS~17
et le capital de solvabilité requis (SCR) sont hors périmètre.
\end{encadre}

Les garanties sont traitées par jalons~:

\begin{itemize}[itemsep=2pt]
\item première version~: décès et arrêt de travail, avec les garanties
  liées\fsp;
\item deuxième version~: épargne en euros et en unités de compte, avec
  l'actif facultatif et la projection actif-passif sous scénarios
  économiques\fsp;
\item troisième version~: retraite, préretraite et indemnités de fin de
  carrière.
\end{itemize}

Le package reçoit une \textbf{base de gestion} (contrats, affiliations,
sinistres, événements, dossiers, prestations, rentes en service,
cotisations) et un \textbf{plan de calcul}, qui déclare toutes les
hypothèses, sans valeur par défaut. Les tables (réglementaires ou
d'expérience) sont toujours des entrées. Un triangle n'est jamais une
entrée~: il est reconstitué de la base sinistres. Le projet ne contient
aucune donnée réelle~: la base brute fictive est tirée par le simulateur du
package, sous une loi vraie et une graine, avec des défauts de données
paramétrés.

% =============================================================================
\section{Conventions}
\label{sec:conventions}
% =============================================================================

\begin{itemize}[itemsep=2pt]
\item \textbf{Vocabulaire.} Les termes du domaine sont ceux du glossaire du
  projet (\code{CONTEXT.md}), employés tels quels.
\item \textbf{Noms.} Les modules, fonctions et colonnes portent des noms
  français sans accents. Une fonction se cite par \code{nom()}, un module
  par son nom, qui est celui de son fichier \code{R/<module>.R}.
\item \textbf{Statuts.} Cinq statuts épistémiques ne se confondent jamais~:
  résultat exact, résultat asymptotique, approximation numérique, constat de
  simulation, jugement d'expert.
\item \textbf{Tableaux de fonctions.} Chaque module a un tableau qui recense
  ses fonctions de premier niveau, hors auxiliaires (fonctions dont le nom
  commence par un point). La batterie de concordance compare ces tableaux
  au code (section~\ref{sec:concordance}).
\item \textbf{Tableaux de constantes.} Les constantes d'un module dont le
  nom est en majuscules sont citées d'après le code, dans un tableau de
  constantes placé dans la partie du module\fsp; la batterie de
  concordance compare au code les valeurs qui y sont écrites sous une
  forme littérale (section~\ref{sec:concordance}).
\item \textbf{Chiffres.} Aucun chiffre issu d'une exécution du package
  n'est écrit dans ce document.
\end{itemize}


\partie{Partie II}{Architecture}

% =============================================================================
\section{Package et point d'entrée}
\label{sec:package}
% =============================================================================

Le moteur est le package R \code{provie}. Il n'emploie que les packages de
priorité \code{base} de R (\code{stats}, \code{tools}, \code{utils}) et ne
dépend d'aucun programme Java~; deux scripts de contrôle le vérifient
(section~\ref{sec:controles}). Il contient toute la logique quantitative.
L'application Shiny est hors du package~: elle n'appelle que ses fonctions
exportées et ne calcule rien.

Le package a un \textbf{point d'entrée unique}, porté par le module du
même nom (section~\ref{sec:calculer-provisions})~:

\begin{center}
\code{calculer\_provisions(donnees, plan, tables = list(), provenance = NULL)}
\end{center}

\noindent Il reçoit la base sous forme
d'une liste de \code{data.frame}, nommés comme dans le dictionnaire de la
base, un plan de calcul et, en arguments facultatifs, les tables importées
et la provenance des fichiers lus~: jamais un fichier. Il exécute le plan
d'un bloc et retourne un objet complet, de classe \code{resultat\_provie},
dans lequel l'application, le rapport figé et les tests lisent sans
recalculer.

Le moteur ne parle pas~: il n'émet ni avertissement, ni message, ni sortie
à la console, et ne lit pas la date du système\fsp; un script de contrôle
le vérifie (section~\ref{sec:controles}). Quand le calcul entier est
indéfini, il lève un
\textbf{refus d'appel}, qui interrompt l'appel sans rendre d'objet et porte
un motif d'une liste fermée (section~\ref{sec:outils}).

Dans l'état du code, le package exporte huit fonctions~:
\code{dictionnaire()} et \code{controler\_types()} (module
\code{dictionnaire}), \code{lire\_base()}, \code{ecrire\_base()},
\code{ecrire\_jeu()} et \code{lire\_jeu()} (module \code{lecture}),
\code{controler\_base()} (module \code{donnees}) et \code{simuler\_base()}
(module \code{simulation}), ainsi que les méthodes \code{format()} et
\code{print()} de la classe \code{jeu\_simule}. Le point d'entrée
\code{calculer\_provisions()} n'est pas encore écrit.

% =============================================================================
\section{Quinze modules et leur graphe}
\label{sec:modules}
% =============================================================================

Le package compte quinze modules, un fichier \code{R/<module>.R} par
module. Un fichier de \code{R/} hors de cette liste est un écart relevé par
la batterie de concordance et par le script de contrôle du graphe des
modules (section~\ref{sec:controles}). Un fichier de module n'existe qu'à partir du
moment où le module a un contenu.

La colonne «~Peut appeler~» est le graphe des appels autorisés entre
modules. Deux modules sont hors de la chaîne du moteur~: \code{lecture}
(lecture et écriture des fichiers) et \code{simulation} (simulateur)\fsp;
aucun module du moteur ne les appelle. Un poste qui combine deux familles de
calcul est une étape de l'orchestration, qui appelle chaque famille et
applique la règle de combinaison déclarée par le catalogue des
configurations de garanties\fsp; la famille «~triangle~» n'appelle pas la
famille «~engagement~». Le script de contrôle du graphe des modules
vérifie, sur le source, que chaque module n'appelle que les modules de sa
colonne «~Peut appeler~» (section~\ref{sec:controles}).

{\small
\renewcommand{\arraystretch}{1.22}
\begin{longtable}{@{}p{3.7cm}>{\raggedright\arraybackslash}p{6.9cm}>{\raggedright\arraybackslash}p{4.5cm}@{}}
\toprule
\textbf{Module} & \textbf{Rôle} & \textbf{Peut appeler} \\
\midrule
\endfirsthead
\toprule
\textbf{Module} & \textbf{Rôle} & \textbf{Peut appeler} \\
\midrule
\endhead
\multicolumn{3}{@{}l@{}}{\textit{Suite page suivante}} \\
\endfoot
\bottomrule
\endlastfoot
\code{outils} & Services sans contenu métier~: refus, graine locale, empreinte, écriture canonique, sérialisation, plateforme, conventions d'âge, constantes & Aucun module \\
\code{dictionnaire} & Lecture et validation du dictionnaire lisible par machine\fsp; accès par table et par colonne & \code{outils} \\
\code{lecture} & Lecture et écriture des fichiers, types imposés par le dictionnaire\fsp; hors de la chaîne du moteur & \code{dictionnaire}, \code{outils} \\
\code{donnees} & Contrôles de la base, traitements déclarés des défauts, vue à une date, segments & \code{dictionnaire}, \code{outils} \\
\code{tables} & Formats des tables, représentation commune, tables embarquées, courbes des taux, plafonds réglementaires & \code{outils} \\
\code{garanties} & Lecture et validation du catalogue des configurations de garanties & \code{dictionnaire}, \code{outils} \\
\code{engagement} & Famille «~engagement~»~: échéanciers, valeur actuelle probable sous le modèle multi-états & \code{tables}, \code{garanties}, \code{donnees}, \code{outils} \\
\code{triangle} & Famille «~triangle~»~: reconstitution, estimateurs, diagnostics & \code{garanties}, \code{donnees}, \code{outils} \\
\code{risque\_en\_cours} & Famille «~risque en cours~»~: limites du contrat, provision pour risques croissants & \code{engagement}, \code{tables}, \code{garanties}, \code{donnees}, \code{outils} \\
\code{regles\_comptables} & Catégorie «~règle comptable~» & \code{garanties}, \code{tables}, \code{donnees}, \code{outils} \\
\code{projection} & Famille «~projection~» (deuxième version) & \code{engagement}, \code{tables}, \code{garanties}, \code{donnees}, \code{outils} \\
\code{plan} & Validation et canonisation du plan, fiche de calcul, sérialisation & \code{tables}, \code{garanties}, \code{dictionnaire}, \code{donnees}, \code{outils} \\
\code{rapport} & Représentation canonique du résultat, empreintes, quantités restituées, tableau de passage, comparaison, rapport figé & \code{outils}, \code{dictionnaire}, \code{donnees}, \code{tables}, \code{garanties}, \code{engagement}, \code{triangle}, \code{risque\_en\_cours}, \code{regles\_comptables}, \code{projection}, \code{plan} \\
\code{simulation} & Simulateur de la base brute fictive\fsp; hors de la chaîne du moteur & \code{dictionnaire}, \code{garanties}, \code{tables}, \code{outils} \\
\code{calculer\_provisions} & Point d'entrée, table des étapes, combinaison des postes à deux familles, rejeu d'une fiche & Tous les modules, sauf \code{lecture} et \code{simulation} \\
\end{longtable}
}


\partie{Partie III}{Module outils}

% =============================================================================
\section{Rôle et fonctions du module outils}
\label{sec:outils}
% =============================================================================

Le module \code{outils} (\code{R/outils.R}) rend des services sans contenu
métier aux autres modules~: refus classés, graine locale, empreinte,
écriture canonique des valeurs, sérialisation d'une liste nommée en texte,
description de la plateforme, calcul de l'âge selon une convention
déclarée. Il n'appelle aucun autre module. Toute source
d'aléa du package passe par \code{avec\_graine()}~: aucun autre fichier de
\code{R/} ne règle le générateur pseudo-aléatoire, ce que vérifie un script
de contrôle (section~\ref{sec:controles}). Aucune fonction du module n'est
exportée.

\begin{center}
\small
\rowcolors{2}{tablerow}{white}
% concordance-fonctions : outils
\begin{tabular}{@{}p{3.9cm}>{\raggedright\arraybackslash}p{3.0cm}p{8.2cm}@{}}
\rowcolor{tableheader}
\textcolor{financeblue}{\textbf{Fonction}} & \textcolor{financeblue}{\textbf{Arguments}} & \textcolor{financeblue}{\textbf{Rôle}} \\
\code{refuser()} & \code{message}, \code{motif}, champs facultatifs & Lève un refus d'appel~: condition de classe \code{refus\_provie} (qui hérite de \code{error}), qui interrompt l'appel sans rendre d'objet\fsp; le motif est pris dans la liste fermée \code{MOTIFS\_REFUS}\fsp; les champs facultatifs (\code{module}, \code{table}, \code{colonne}, \code{usage}, \code{nombre}, \code{identifiants}, \code{attendu}, \code{recu}) situent le refus, les identifiants étant tronqués à dix. \\
\code{avec\_graine()} & \code{graine}, \code{expr} & Évalue \code{expr} sous le générateur Mersenne-Twister (échantillonneurs Inversion et Rejection) et la graine donnée, puis rétablit les genres du générateur et la graine de l'appelant, même en cas d'erreur (section~\ref{sec:graine}). \\
\code{empreinte()} & \code{octets} & Rend le MD5 d'un vecteur d'octets (\code{raw}), en 32 caractères hexadécimaux minuscules, calculé par \code{tools::md5sum()} sur un fichier temporaire supprimé à la sortie (R~4.3.3 ne calcule le MD5 que d'un fichier). \\
\code{ecrire\_valeurs()} & \code{x}, \code{forme = "texte"} & Écrit un vecteur atomique simple (logique, entier, réel, caractère, date) sous sa forme canonique~: en texte, une chaîne par valeur (réels à 17 chiffres significatifs, dates ISO~8601)\fsp; en octets, sans étiquette de type, nombres en petit-boutiste (section~\ref{sec:ecriture}). \\
\code{ecrire\_liste()} & \code{x} & Sérialise une liste nommée en texte aplati UTF-8, une ligne par feuille (chemin, type, valeurs), rendu en octets\fsp; n'écrit aucun fichier (section~\ref{sec:format-liste}). \\
\code{lire\_liste()} & \code{octets} & Relit, sans évaluer de code, les octets écrits par \code{ecrire\_liste()}\fsp; refuse tout texte qui n'est pas canonique, de sorte que l'aller-retour rend une liste identique. \\
\code{decrire\_plateforme()} & Aucun & Décrit la plateforme d'exécution pour la fiche de calcul (version de R, système, machine, boutisme, bibliothèques d'algèbre linéaire, d'expressions régulières et Unicode), en liste sérialisable par \code{ecrire\_liste()}, sans date ni nom d'utilisateur\fsp; chaque champ est une chaîne de longueur~1, et un champ indisponible (\code{Sys.info()} qui rend \code{NULL}, composant absent de \code{Sys.info()} ou d'\code{extSoftVersion()}) vaut \code{NA}. \\
\code{calculer\_age()} & \code{date\_naissance}, \code{date}, \code{convention} & Rend l'âge entier, en années, à une date, selon une convention de la liste fermée \code{CONVENTIONS\_AGE} (\code{revolu}, \code{millesime}, \code{plus\_proche}, \code{naissance\_31\_decembre})\fsp; né un 29~février, l'anniversaire d'une année non bissextile est le 1er~mars\fsp; \code{NA} si la date précède la naissance\fsp; refuse par \code{plan\_invalide} une convention hors de la liste. Seule fonction d'âge du package. \\
\code{date\_atteinte\_age()} & \code{date\_naissance}, \code{age}, \code{convention} & Fonction inverse de \code{calculer\_age()}~: première date, au plus tôt la naissance, à laquelle l'âge selon la convention vaut \code{age}. \\
\end{tabular}
\end{center}

Le module porte aussi sept constantes (section~\ref{sec:outils-constantes}).
Les deux fonctions d'âge sont décrites à la section~\ref{sec:age}.

% =============================================================================
\section{Constantes et motifs de refus}
\label{sec:outils-constantes}
% =============================================================================

Les constantes sont citées d'après le code (\code{R/outils.R}).

\begin{center}
\small
\rowcolors{2}{tablerow}{white}
\begin{tabular}{@{}p{5.6cm}p{3.7cm}p{5.8cm}@{}}
\rowcolor{tableheader}
\textcolor{financeblue}{\textbf{Constante}} & \textcolor{financeblue}{\textbf{Valeur}} & \textcolor{financeblue}{\textbf{Objet}} \\
\code{VERSION\_FICHE} & entier \code{1} & Version du format de la fiche de calcul. Aucune fonction ne la lit encore. \\
\code{VERSION\_DICTIONNAIRE\_ATTENDUE} & \code{"1.5"} & Version du dictionnaire de la base attendue par le moteur (dictionnaire de la base, version~1.5). Lue par \code{dictionnaire()}, qui refuse un dictionnaire embarqué d'une autre version, et par \code{controler\_version\_extraction()}, qui refuse une base d'une autre version (section~\ref{sec:dictionnaire}). \\
\code{REPRESENTATION\_CANONIQUE} & \code{"provie-resultat-1"} & Identifiant de la représentation canonique du résultat. Aucune fonction ne la lit encore. \\
\code{VERSION\_LISTE} & entier \code{1} & Version du format de sérialisation des listes (section~\ref{sec:format-liste}). \\
\code{MOTIFS\_REFUS} & dix-neuf motifs & Liste fermée des motifs de refus d'appel (ci-dessous). \\
\code{CONVENTIONS\_AGE} & quatre chaînes & Liste fermée des conventions de calcul de l'âge~: \code{revolu}, \code{millesime}, \code{plus\_proche}, \code{naissance\_31\_decembre} (section~\ref{sec:age}). \\
\code{AGE\_MAXIMUM} & entier \code{200} & Âge le plus élevé admis par \code{date\_atteinte\_age()}~; au-delà, la fonction lève une erreur R ordinaire (section~\ref{sec:age}). \\
\end{tabular}
\end{center}

\textbf{Refus d'appel.} La fonction \code{refuser()} lève une condition R
de classe \code{refus\_provie}, qui hérite de \code{error}~: l'appel est
interrompu sans rendre d'objet, et l'appelant peut attraper la condition
par sa classe et lire son motif. Le motif doit appartenir à la liste fermée
\code{MOTIFS\_REFUS}. Un motif hors de la liste, ou un champ facultatif mal
formé (\code{module}, \code{table}, \code{colonne} ou \code{usage} qui n'est
pas une chaîne unique, \code{nombre} qui n'est ni une valeur numérique
unique entière et finie, au plus égale en valeur absolue au plus grand
entier de R, ni une valeur manquante \code{NA} unique d'un type atomique
quelconque, \code{NaN} étant refusé, \code{attendu} ou \code{recu} qui n'est pas une valeur atomique unique), est
une erreur de programmation~: \code{refuser()} lève alors une erreur R
ordinaire, qui n'est pas un refus d'appel. Les \code{identifiants} sont
tronqués aux dix premiers\fsp; \code{nombre}, converti en entier, porte
le total quand l'appelant le renseigne.

La liste fermée compte dix-neuf motifs. Le tableau ci-dessous dit, pour
chacun, quelles fonctions le lèvent dans l'état du code, directement ou par
une fonction qu'elles appellent (ainsi, toute fonction qui lit le
dictionnaire peut lever \code{dictionnaire\_invalide})\fsp; un motif levé
par aucune fonction est réservé aux modules qui le lèveront, et son sens se
documentera avec eux.

\begin{center}
\small
\rowcolors{2}{tablerow}{white}
\begin{tabular}{@{}p{5.6cm}p{9.5cm}@{}}
\rowcolor{tableheader}
\textcolor{financeblue}{\textbf{Motif}} & \textcolor{financeblue}{\textbf{Levé, dans l'état du code, par}} \\
\code{plan\_incomplet} & \code{controler\_base()}~: date d'arrêté absente. \\
\code{plan\_invalide} & \code{controler\_base()}~: date d'arrêté mal formée, déclaration ou convention mal formée ou hors de sa liste fermée\fsp; \code{calculer\_age()} et \code{date\_atteinte\_age()}~: convention hors de \code{CONVENTIONS\_AGE}. \\
\code{liste\_invalide} & \code{ecrire\_valeurs()}, \code{ecrire\_liste()} et \code{lire\_liste()}~: valeur, liste ou texte que l'écriture canonique ou le format de liste n'admet pas (sections~\ref{sec:ecriture} et~\ref{sec:format-liste})\fsp; \code{controler\_types()}, \code{constater\_forme()} et \code{controler\_base()}~: base qui n'est pas une liste nommée de \code{data.frame} aux tables connues\fsp; \code{ecrire\_base()}~: valeur que le CSV ne relirait pas à l'identique. \\
\code{fichier\_invalide} & \code{lire\_base()}, \code{ecrire\_base()}, \code{ecrire\_jeu()} et \code{lire\_jeu()}~: chemin qui n'est pas une chaîne unique, dossier introuvable, non vide ou impossible à créer, fichier inconnu du dictionnaire, fichier dont la lecture ne rendrait pas le texte exact (section~\ref{sec:lecture-base})\fsp; pour \code{lire\_jeu()}, disposition d'un jeu, liste d'empreintes, volumétrie ou tables de vérité non conformes (section~\ref{sec:lecture-jeu}). \\
\code{table\_base\_absente} & \code{controler\_base()}~: table obligatoire absente. \\
\code{colonne\_absente} & \code{controler\_types()}, et par elle \code{ecrire\_base()}~: colonne obligatoire ou conditionnelle absente\fsp; \code{lire\_base()} et \code{controler\_base()}, pour la seule table \code{extraction}. \\
\code{colonne\_inconnue} & Mêmes fonctions que \code{colonne\_absente}~: colonne inconnue du dictionnaire. \\
\code{type\_incompatible} & Mêmes fonctions que \code{colonne\_absente}~: colonne d'un autre type que son type déclaré. \\
\code{extraction\_invalide} & \code{controler\_version\_extraction()}, et par elle \code{controler\_types()}, \code{ecrire\_base()}, \code{lire\_base()} et \code{controler\_base()}~: version du dictionnaire autre que \code{VERSION\_DICTIONNAIRE\_ATTENDUE}\fsp; \code{controler\_types()} et \code{controler\_base()}~: table \code{extraction} qui n'a pas exactement une ligne, et, pour \code{controler\_base()} seule, valeur manquante ou hors domaine, date antérieure au 1er~janvier~1901 ou début d'historique qui ne précède pas la date d'extraction. \\
\code{arrete\_posterieur\_extraction} & \code{controler\_base()}~: date d'arrêté postérieure à la date d'extraction. \\
\code{declaration\_manquante} & Aucune fonction. \\
\code{table\_absente} & Aucune fonction. \\
\code{empreinte\_discordante} & \code{lire\_jeu()}~: fichier d'un jeu dont le MD5 diffère de celui de \code{empreintes.md5}, avec l'empreinte attendue et l'empreinte reçue. \\
\code{version\_discordante} & \code{lire\_liste()}~: en-tête d'une autre version du format de liste, avec la version attendue (\code{attendu}) et la version reçue (\code{recu}). \\
\code{lien\_non\_declare} & Aucune fonction. \\
\code{catalogue\_invalide} & Aucune fonction. \\
\code{dictionnaire\_invalide} & \code{dictionnaire()}, et par elle toute fonction qui lit le dictionnaire~: dictionnaire embarqué incohérent (par exemple d'une autre version, d'une empreinte discordante ou d'une référence interne invalide). \\
\code{format\_table\_inconnu} & Aucune fonction. \\
\code{jeu\_invalide} & \code{simuler\_base()}~: argument qui ne définit pas un jeu simulé (loi vraie, graines, date d'extraction, nom, description, date de visa, déclarations) ou contraire aux contrôles d'entrée (section~\ref{sec:simulation-tirage})\fsp; \code{ecrire\_jeu()}~: objet qui n'a pas la forme d'un jeu simulé ou que \code{lire\_jeu()} ne relirait pas identique. \\
\end{tabular}
\end{center}

% =============================================================================
\section{Graine locale}
\label{sec:graine}
% =============================================================================

\code{avec\_graine(graine, expr)} est la seule voie vers l'aléa dans le
package. La graine doit être une valeur numérique entière et finie, au plus
égale en valeur absolue au plus grand entier de R (\code{.Machine\$integer.max})\fsp;
sinon la fonction lève une erreur R ordinaire. La fonction procède ainsi~:

\begin{enumerate}[itemsep=2pt]
\item elle relève l'état global de l'appelant~: la présence et la valeur de
  la graine courante (\code{.Random.seed}), et les trois genres du
  générateur (\code{RNGkind()})~: générateur uniforme, échantillonneur de
  la loi normale, échantillonneur de \code{sample()}\fsp;
\item elle fixe les genres Mersenne-Twister, Inversion et Rejection, puis la
  graine (\code{set.seed()})\fsp;
\item elle évalue \code{expr}, dans l'environnement de l'appelant, et rend
  sa valeur\fsp;
\item à la sortie, que l'évaluation ait réussi ou non, elle rétablit les
  trois genres de l'appelant, puis sa graine~: la graine est réaffectée si
  elle existait, supprimée sinon. Le rétablissement de la graine est fait
  même si celui des genres échoue.
\end{enumerate}

Rétablir l'échantillonneur «~Rounding~» de \code{sample()} (celui de R
antérieur à la version~3.6.0) fait émettre par R un avertissement. Quand
l'échantillonneur de l'appelant est «~Rounding~», le rétablissement des
genres est fait sous \code{suppressWarnings()}, par décision du
mainteneur du 7~novembre~2026 (motif ci-dessous). Dans tous les autres cas, un avertissement
émis par le rétablissement remonte à l'appelant, et devient une erreur si
les avertissements sont convertis en erreurs (\code{options(warn = 2)},
comme dans les batteries de tests)~; c'est le cas, par décision du
mainteneur du même jour, de l'avertissement de l'échantillonneur de la loi
normale «~Buggy Kinderman-Ramage~». Le motif de cette différence, donné par
le mainteneur le 8~novembre~2026~: «~Rounding~» est l'état ordinaire d'un
appelant qui rejoue d'anciens résultats par \code{RNGversion()}, et son
avertissement n'a pas à sortir d'une restauration\fsp; «~Buggy
Kinderman-Ramage~» est un générateur défectueux, dont l'usage doit être
signalé. Les genres et la graine de l'appelant sont rétablis dans
les deux cas. Une limite subsiste~: si l'appelant combine «~Buggy
Kinderman-Ramage~» et «~Rounding~», l'avertissement du premier est étouffé
avec celui du second (issue~27 du dépôt).

Statut~: la batterie de reproductibilité constate, sur la plateforme où
elle tourne, que deux appels de même graine rendent les mêmes tirages, que
la graine et les genres de l'appelant sont rétablis à l'identique, y
compris après une erreur dans \code{expr}, et que la suite des tirages de
l'appelant est inchangée. Elle compare aussi trois tirages sous une graine
fixée à des valeurs écrites en dur, à une tolérance relative déclarée~; en
intégration continue, cette comparaison tourne sur la plateforme de
référence (section~\ref{sec:batteries}). La fonction ne garantit pas, à
elle seule, l'identité des tirages d'une plateforme à l'autre.

% =============================================================================
\section{Écriture canonique des valeurs}
\label{sec:ecriture}
% =============================================================================

\code{ecrire\_valeurs(x, forme)} écrit un vecteur atomique simple~: logique,
entier, réel ou caractère sans aucun attribut, ou date (vecteur de réels qui
porte le seul attribut de classe \code{Date}). Tout autre objet (facteur,
vecteur nommé, complexe, liste\ldots) est refusé (\code{liste\_invalide}).

\textbf{Forme texte} (\code{forme = "texte"}), une chaîne par valeur~:

\begin{itemize}[itemsep=2pt]
\item logique~: \code{T}, \code{F}, \code{NA}\fsp;
\item entier~: écriture décimale, \code{NA}\fsp;
\item réel~: 17 chiffres significatifs (format \code{\%.17g}), jetons
  \code{NA}, \code{NaN}, \code{Inf} et \code{-Inf}\fsp; le zéro négatif
  s'écrit \code{0}. Dix-sept chiffres significatifs suffisent à retrouver
  tout réel double précision (binaire64 de la norme IEEE~754) lorsque les
  deux conversions, vers le décimal et depuis le décimal, sont
  correctement arrondies. Dans \code{provie}, l'écriture passe par
  \code{sprintf()}, donc par la bibliothèque C\fsp; la relecture passe par
  \code{as.numeric()}, c'est-à-dire par le convertisseur propre de R
  (\code{R\_strtod}), qui accumule les chiffres et les puissances de dix en
  \code{long double}~: il n'est pas garanti correctement arrondi, et sa
  précision dépend du \code{long double} de la plateforme. L'aller-retour
  exact est donc un constat~: la batterie de reproductibilité le vérifie
  sur trois mille réels, relus au bit près, sur la plateforme où elle
  tourne\fsp;
\item date~: ISO~8601, \code{AAAA-MM-JJ}, l'année complétée à quatre chiffres
  (l'an~999 s'écrit \code{0999})\fsp; \code{NA}. Une date non entière, \code{NaN},
  infinie ou dont l'année ne tient pas sur quatre chiffres est refusée
  (\code{liste\_invalide})\fsp;
\item caractère~: la chaîne en UTF-8, la barre oblique inverse, la
  tabulation, le saut de ligne et le retour chariot étant échappés en
  \code{\textbackslash\textbackslash}, \code{\textbackslash t},
  \code{\textbackslash n} et \code{\textbackslash r}, les autres caractères
  de contrôle et le caractère DEL en \code{\textbackslash x} suivi de deux
  chiffres hexadécimaux minuscules\fsp; une valeur manquante s'écrit
  \code{\textbackslash N}.
\end{itemize}

\textbf{Forme octets} (\code{forme = "octets"}), un vecteur \code{raw} sans
étiquette de type (le type doit donc être connu de celui qui relit)~:
logique sur un octet (0, 1, et 2 pour \code{NA})\fsp; entier sur quatre
octets petit-boutistes\fsp; réel sur huit octets IEEE~754 petit-boutistes,
\code{NA} et \code{NaN} sur deux motifs fixes (\code{7FF00000000007A2} et
\code{7FF8000000000000} en hexadécimal, octet de poids fort en tête), le
zéro négatif écrit comme \code{0}\fsp; caractère par sa longueur en octets
UTF-8 sur quatre octets petit-boutistes (\code{FFFFFFFF} pour \code{NA}),
suivie de ses octets UTF-8\fsp; date comme sa chaîne ISO~8601.

\textbf{Encodage des chaînes.} Les chaînes sont ramenées en UTF-8 sans
jamais être corrigées~: une chaîne ASCII est prise telle quelle\fsp; une
chaîne marquée UTF-8 est validée\fsp; une chaîne marquée Latin-1 est
traduite en UTF-8\fsp; une chaîne marquée «~bytes~» est refusée. Une chaîne
non ASCII sans marque d'encodage est refusée si la locale de la session
n'est pas UTF-8, et validée comme UTF-8 sinon. Le refus vient de ce que R
ne sait pas dire, hors locale UTF-8, quel est l'encodage d'une chaîne non
marquée~: la convertir écrirait d'autres octets que les siens (constat du
7~novembre~2026 sous R~4.3.3~: les deux octets de «~é~» sont devenus le
texte \code{<c3><a9>}). Une chaîne qui n'est pas de l'UTF-8 valide est
refusée. Tous ces refus portent le motif \code{liste\_invalide}. Cette règle
vaut aussi pour les noms d'une liste sérialisée (section~\ref{sec:format-liste})\fsp;
une chaîne lue d'un fichier doit donc porter la marque UTF-8 pour être
écrite sous toute locale.

% =============================================================================
\section{Format provie-liste 1}
\label{sec:format-liste}
% =============================================================================

\code{ecrire\_liste(x)} sérialise une liste nommée en un texte UTF-8, rendu
en octets (l'écriture d'un fichier revient à l'appelant). Le format porte le
numéro \code{VERSION\_LISTE}~:

\begin{itemize}[itemsep=2pt]
\item lignes terminées par un saut de ligne, y compris la dernière\fsp;
  première ligne \code{provie-liste 1}\fsp;
\item puis une ligne par feuille, dans l'ordre des éléments de la liste,
  sous-listes parcourues en profondeur, faite de champs séparés par une
  tabulation~: le chemin, le type, puis une valeur par élément du
  vecteur (aucune pour un vecteur vide), écrite sous la forme texte de la
  section~\ref{sec:ecriture}\fsp;
\item le chemin joint par \code{/} les noms des ancêtres et de la feuille,
  chaque nom échappé comme une chaîne, et le caractère \code{/} d'un nom
  écrit \code{\textbackslash x2f}\fsp;
\item le type est l'un de \code{null}, \code{liste} (liste vide),
  \code{logical}, \code{integer}, \code{double}, \code{character},
  \code{date}.
\end{itemize}

Une feuille est \code{NULL}, une liste vide ou un vecteur atomique simple.
Les éléments d'une liste non vide ont des noms renseignés et distincts, et
aucune liste ne porte d'autre attribut que ses noms\fsp; sinon l'écriture est
refusée (\code{liste\_invalide}), avec un message propre à chacune des deux
causes (noms absents ou non distincts, attribut non admis). Une liste vide s'écrit par la seule ligne
d'en-tête. Par exemple, la liste \code{list(a = 1.5, b = list(c = "x", d =
NULL))} s'écrit en quatre lignes, en-tête compris (la tabulation est figurée ici par un
espacement)~:

\begin{encadre}
\small
\code{provie-liste 1}\\
\code{a}\quad\code{double}\quad\code{1.5}\\
\code{b/c}\quad\code{character}\quad\code{x}\\
\code{b/d}\quad\code{null}
\end{encadre}

\textbf{Lecture stricte.} La fonction \code{lire\_liste(octets)} relit ce
texte sans évaluer de code. Elle refuse, par le motif
\code{liste\_invalide}, un contenu qui contient un octet nul, qui n'est pas de l'UTF-8 valide ou qui ne finit pas
par un saut de ligne, un en-tête absent, une ligne sans chemin ou sans
type, un chemin mal formé, un chemin dupliqué ou une feuille placée sous
une autre (comparaison faite sur les noms décodés~: \code{a} et
\code{\textbackslash x61} sont le même nom), un jeton qui n'est pas de son type, une date inexistante, une
séquence d'échappement inconnue ou l'octet nul échappé. Un en-tête
\code{provie-liste} d'une autre version est refusé par le motif
\code{version\_discordante}. La liste est ensuite reconstruite niveau par
niveau~: les chemins sont regroupés selon leur premier nom, dans l'ordre de
première apparition, et chaque groupe est construit de la même façon au
niveau suivant. Le travail est ainsi proportionnel au nombre de feuilles
multiplié par la profondeur, là où une insertion feuille par feuille
recopierait le nœud parent à chaque feuille (issue~25 du dépôt). Enfin, la liste relue est réécrite, et la
lecture est refusée si la réécriture diffère, ne serait-ce que d'un octet,
du texte reçu (\code{1.50} au lieu de \code{1.5}, par exemple)~: seul le
texte canonique est accepté, par décision du mainteneur du
7~novembre~2026. Il en résulte, par construction, que tout texte accepté
est restitué à l'octet près par \code{ecrire\_liste(lire\_liste(octets))}.
Que \code{lire\_liste(ecrire\_liste(x))} rende une liste identique à
\code{x} est un constat (batterie de reproductibilité, sur une liste qui
couvre les sept types, les valeurs manquantes, les réels spéciaux, les
chaînes échappées et la description de la plateforme), non une
démonstration. «~Identique~» s'entend au sens d'\code{identical()} avec
ses réglages par défaut, qui ne distingue pas \code{-0} de \code{0} et
ignore la marque d'encodage des chaînes~: un zéro négatif est relu
\code{0} (section~\ref{sec:ecriture}), et une chaîne marquée Latin-1 est
relue marquée UTF-8.

% =============================================================================
\section{Conventions d'âge}
\label{sec:age}
% =============================================================================

\code{calculer\_age(date\_naissance, date, convention)} est la seule
fonction de calcul de l'âge du package~: le module \code{donnees} l'emploie
pour l'âge invraisemblable (défaut D02c) et pour la fin de charge d'un
enfant (défaut D07, sous-cas~(c)), et le simulateur emploie sa fonction
inverse. Elle est placée dans \code{outils} parce qu'elle ne fait que de
l'arithmétique de calendrier~: le choix métier, la convention, reste une
déclaration du plan (\code{conventions\$age}), sans valeur par défaut. Elle
rend un âge entier, en années, selon l'une des quatre conventions de la
liste fermée \code{CONVENTIONS\_AGE}~:

\begin{itemize}[itemsep=2pt]
\item \code{revolu}~: nombre d'anniversaires atteints au plus tard à la
  date\fsp;
\item \code{millesime}~: année de la date moins année de naissance\fsp;
\item \code{plus\_proche}~: âge révolu, augmenté d'un an si l'anniversaire
  suivant est plus proche de la date, en jours, que le précédent\fsp; à
  égalité, l'âge supérieur\fsp;
\item \code{naissance\_31\_decembre}~: âge révolu calculé sur la naissance
  reportée au 31~décembre le plus voisin, celui de l'année précédente si la
  naissance tombe au plus tard le 1er~juillet (1er~juillet d'une année
  bissextile compris), celui de l'année de naissance sinon\fsp; entre la
  naissance et un 31~décembre reporté qui lui est postérieur, l'âge vaut~0.
\end{itemize}

Pour une personne née un 29~février, l'anniversaire d'une année non
bissextile est le 1er~mars. L'âge est manquant (\code{NA}) si l'une des
dates manque ou si la date précède la naissance. Une convention hors de la
liste est refusée par le motif \code{plan\_invalide}\fsp; une date qui
n'est pas de classe \code{Date}, ou qui n'est pas entière et finie, est une
erreur de programmation (erreur R ordinaire). Les deux vecteurs de dates
ont la même longueur, ou l'un est de longueur~1.

\code{date\_atteinte\_age(date\_naissance, age, convention)} est la
fonction inverse~: elle rend la première date, au plus tôt la naissance, à
laquelle l'âge selon la convention vaut \code{age}, de sorte que
\code{calculer\_age(n, date\_atteinte\_age(n, a, c), c)} vaut~\code{a}.
L'âge demandé est un entier de~0 à \code{AGE\_MAXIMUM}, ou \code{NA}\fsp;
au-delà, la fonction lève une erreur R ordinaire plutôt que de rendre une
valeur manquante sans le dire.

Statuts~: les quatre définitions sont celles de la spécification du module
\code{donnees} visée par le mainteneur (décisions M151 et M154)\fsp; la convention
\code{naissance\_31\_decembre} est retenue \textbf{par analogie}, jugement
d'expert, et n'est pas présentée comme exigée par un texte pour les
contrôles de la base~; la spécification la déconseille pour D02c. Pour la fin de charge d'un enfant (D07~(c)), seules
\code{revolu} et \code{millesime} sont admises, une clause de limite d'âge
se lisant sur la date de naissance réelle (section~\ref{sec:donnees-bilan}).
L'aller-retour entre les deux fonctions est vérifié par la batterie
unitaire\fsp; c'est un constat sur les cas testés, non une démonstration.


\partie{Partie IV}{Module dictionnaire}

% =============================================================================
\section{Rôle et fonctions du module dictionnaire}
\label{sec:dictionnaire}
% =============================================================================

Le module \code{dictionnaire} lit et valide le dictionnaire lisible par
machine, transcription du dictionnaire de la base, et donne accès, par table
et par colonne, aux types, domaines, obligations, modes de saisie et défauts
de données. Il porte aussi l'automate des événements, la liste des
paramètres, la liste des défauts et celle des codes de contrôle de
structure. Il ne peut appeler que \code{outils}. Deux de ses cinq
fonctions sont exportées, \code{dictionnaire()} et
\code{controler\_types()}\fsp; les trois autres servent les modules
\code{lecture} et \code{donnees}.

\begin{center}
\small
\rowcolors{2}{tablerow}{white}
% concordance-fonctions : dictionnaire
\begin{tabular}{@{}p{5.6cm}>{\raggedright\arraybackslash}p{3.0cm}p{6.5cm}@{}}
\rowcolor{tableheader}
\textcolor{financeblue}{\textbf{Fonction}} & \textcolor{financeblue}{\textbf{Arguments}} & \textcolor{financeblue}{\textbf{Rôle}} \\
\code{dictionnaire()} & Aucun & Lit et valide, à chaque appel, le dictionnaire lisible par machine embarqué sous \code{inst/extdata/dictionnaire/} (sept CSV et \code{dictionnaire.txt}), et le rend en liste nommée~: version, date de visa, puis un \code{data.frame} par CSV, dont \code{controles}\fsp; refuse par \code{dictionnaire\_invalide} un dictionnaire incohérent (section~\ref{sec:dictionnaire-csv}). \\
\code{controler\_types()} & \code{donnees} & Contrôle, sans rien convertir, en trois temps~: la structure de la liste, la table \code{extraction} (version, une ligne), puis les écarts de forme de toutes les tables, tous relevés et comptés avant le refus\fsp; rend la base inchangée ou la refuse (section~\ref{sec:dictionnaire-csv}). \\
\code{constater\_forme()} & \code{donnee}, \code{table}, \code{dico} & Relève, sans s'interrompre ni rien convertir, les écarts de forme d'une table (colonne inconnue, colonne obligatoire ou conditionnelle absente, type R) et les rend en \code{data.frame}, une ligne par écart, avec son code de \code{controles.csv} et les lignes concernées\fsp; non exportée, appelée par \code{controler\_types()} et \code{controler\_base()}. \\
\code{lire\_texte\_type()} & \code{x}, \code{type} & Lit une colonne de texte dans un type du dictionnaire par la grammaire stricte, sans convertir ce qui n'est pas exactement l'écriture du type\fsp; rend la colonne typée si toutes ses valeurs se lisent, et les numéros des valeurs illisibles\fsp; non exportée, appelée par \code{lire\_base()} et \code{constater\_forme()}. \\
\code{controler\_version\_extraction()} & \code{donnee} & Refuse (\code{extraction\_invalide}) une table \code{extraction} d'une autre version du dictionnaire que celle qu'attend le moteur\fsp; non exportée, appelée par \code{controler\_types()}, \code{lire\_base()} et \code{controler\_base()}. \\
\end{tabular}
\end{center}

% =============================================================================
\section{Dictionnaire lisible par machine}
\label{sec:dictionnaire-csv}
% =============================================================================

Le dictionnaire lisible par machine est la transcription, sous
\code{inst/extdata/dictionnaire/}, du dictionnaire de la base visé par le
mainteneur (\code{docs/specifications/dictionnaire-base-v1.md}, version~1.5,
visée le 8~novembre~2026). Un script de contrôle du dépôt compare la
transcription au texte (section~\ref{sec:controles}). Il compte sept
fichiers CSV et un en-tête~:

\begin{center}
\small
\rowcolors{2}{tablerow}{white}
\begin{tabular}{@{}p{3.9cm}p{11.2cm}@{}}
\rowcolor{tableheader}
\textcolor{financeblue}{\textbf{Fichier}} & \textcolor{financeblue}{\textbf{Contenu}} \\
\code{dictionnaire.txt} & Version du dictionnaire, date de visa et empreinte MD5 de chacun des sept CSV, au format \code{provie-liste 1} (section~\ref{sec:format-liste}). \\
\code{tables.csv} & Une ligne par table de la base~: numéro, nom, objet, clé primaire, obligation (\code{O}, \code{F}, ou \code{F\_si:<table>}), visibilité. \\
\code{colonnes.csv} & Une ligne par colonne~: table, nom, type, modalités, domaine, unité, obligation (\code{O}, \code{C}, \code{F}) et condition, mode de saisie et colonne de saisie, date future admise (\code{futur\_admis}), référence, défauts de données qui la concernent. \\
\code{types\_evenements.csv} & Une ligne par type d'événement, dans l'ordre~: niveau, garanties, effet, colonnes requises, caractère programmable, automate du sinistre d'arrêt de travail (état de départ et d'arrivée), motifs. \\
\code{parametres.csv} & Paramètres de garantie~: code de garantie, paramètre, unité, dimension, obligation et condition. \\
\code{defauts.csv} & Défauts de données (codes \code{D..})~: libellé, tables et colonnes, détection (\code{exacte}, \code{presomptive}), traitements admis (\code{R}, \code{X}, \code{H}, \code{S}), paramètre d'intensité et sous-paramètres. \\
\code{filiation.csv} & Correspondance avec la version~1 de la base~: table et colonne, statut (\code{reprise}, \code{renommee}, \code{non\_reprise}), motif. \\
\code{controles.csv} & Codes de contrôle de structure (ci-dessous)~: code, objet, détection, portée, motif de refus d'appel. \\
\end{tabular}
\end{center}

\textbf{Champs.} Un champ à plusieurs valeurs les joint par «~;~». La
colonne \code{futur\_admis} de \code{colonnes.csv} prend l'une des trois
valeurs de \code{FUTUR\_ADMIS\_VALEURS}~: \code{TRUE} (date future
admise), \code{FALSE} (la date d'extraction est une borne haute) et
\code{programmable} (date future admise pour un type d'événement
programmable seulement, valeur réservée à la table \code{evenements}). La
colonne \code{reference} désigne une clé étrangère, simple
(\code{table.colonne}) ou \textbf{composite}
(\code{table.c1;table.c2}, pour la colonne \code{college} de la table
\code{affiliations}). Une référence composite vise une seule table, autre
que celle de la ligne\fsp; elle contient la colonne de la ligne, ne cite
aucune colonne deux fois, et chacune de ses colonnes existe aussi dans la
table de la ligne, sous le même nom. C'est une contrainte d'existence du
n-uplet dans la table visée, qui n'en fait pas une clé de cette table. Une
référence simple ne vise pas la colonne qui la porte.

\textbf{Codes de contrôle de structure.} \code{controles.csv} compte dix
codes, tous à détection exacte~: \code{colonne\_absente},
\code{colonne\_inconnue}, \code{type\_incompatible},
\code{manquant\_obligatoire}, \code{hors\_domaine}, \code{chaine\_vide},
\code{chevauchement}, \code{table\_conditionnelle\_absente},
\code{refus\_induit} et \code{somme\_parts}. Leur portée prend l'une des
quatre valeurs de \code{PORTEES\_CONTROLES}~: \code{colonne}, \code{table},
\code{colonne\_table} (la colonne, ou la table pour une colonne de clé
primaire ou qui porte une référence) et \code{ligne} (signalement). Les
trois premiers codes sont aussi des motifs de refus d'appel, qu'ils portent
dans la colonne \code{motif\_refus\_appel}~: un code qui est un motif de
\code{MOTIFS\_REFUS} porte ce motif, et lui seul.

\textbf{Validation.} \code{dictionnaire()} refuse
(\code{dictionnaire\_invalide}) un en-tête qui ne porte pas, dans cet ordre,
la version, la date de visa et les sept empreintes\fsp; une version autre
que celle qu'attend le moteur (\code{VERSION\_DICTIONNAIRE\_ATTENDUE})\fsp;
un CSV absent, d'une empreinte
discordante, qui n'est pas de l'UTF-8 ou dont les colonnes ne sont pas
celles attendues, et toute référence interne invalide~: table, colonne,
type, défaut, motif ou paramètre cité qui n'existe pas où il est cité,
valeur hors de sa liste fermée, référence mal formée, code de contrôle
répété ou sans objet.

\textbf{Contrôle des types d'une base.} \code{controler\_types(donnees)}
est la porte stricte~: elle ne convertit rien et lève un refus d'appel au
premier des trois temps qui échoue. (1)~Structure (\code{liste\_invalide})~:
liste nommée de \code{data.frame}, aux noms de tables connus, renseignés et
distincts, et aux noms de colonnes renseignés et distincts. (2)~Table
\code{extraction}~: version du dictionnaire, puis une ligne exactement
(\code{extraction\_invalide}). (3)~Écarts de forme de toutes les tables
présentes, relevés par \code{constater\_forme()}~: colonne inconnue,
colonne obligatoire ou conditionnelle absente, type R autre que le type
déclaré. Tous les écarts sont relevés avant le refus, nommés dans le
message et comptés dans le champ \code{nombre}\fsp; le motif
(\code{colonne\_inconnue}, \code{colonne\_absente} ou
\code{type\_incompatible}) est celui du premier écart. Une table
obligatoire absente n'est pas refusée ici~: c'est le contrôle de la base
entière (section~\ref{sec:donnees}).

Les constantes du module sont citées d'après le code
(\code{R/dictionnaire.R})~:

\begin{center}
\small
\rowcolors{2}{tablerow}{white}
\begin{tabular}{@{}p{5.6cm}p{9.5cm}@{}}
\rowcolor{tableheader}
\textcolor{financeblue}{\textbf{Constante}} & \textcolor{financeblue}{\textbf{Valeur et objet}} \\
\code{TYPES\_DICTIONNAIRE} & \code{character}, \code{integer}, \code{numeric}, \code{Date}, \code{logical}~: les cinq types d'une colonne. \\
\code{FUTUR\_ADMIS\_VALEURS} & \code{TRUE}, \code{FALSE}, \code{programmable}~: valeurs de \code{colonnes.csv}, \code{futur\_admis}. \\
\code{PORTEES\_CONTROLES} & \code{colonne}, \code{table}, \code{colonne\_table}, \code{ligne}~: valeurs de \code{controles.csv}, \code{portee}. \\
\code{FICHIERS\_DICTIONNAIRE} & Les sept CSV, dans l'ordre du tableau ci-dessus, qui est celui des empreintes de \code{dictionnaire.txt}. \\
\code{COLONNES\_DICTIONNAIRE} & Liste, par CSV, des noms de colonnes attendus, dans leur ordre. \\
\end{tabular}
\end{center}


\partie{Partie V}{Module lecture}

% =============================================================================
\section{Rôle et fonctions du module lecture}
\label{sec:lecture}
% =============================================================================

Le module \code{lecture} est la couche mince de lecture et d'écriture des
fichiers~: base en CSV, jeux versionnés, tables importées. Les types sont
imposés par le dictionnaire de la base\fsp; une valeur n'est jamais
convertie dans un autre type que le sien. Le module transmet des
\code{data.frame}. Dans l'état du code, il lit et écrit la base en CSV
(section~\ref{sec:lecture-base}) et les jeux simulés
(section~\ref{sec:lecture-jeu})\fsp; la lecture des
tables importées et la provenance des fichiers lus ne sont pas encore
écrites. Il est hors de la chaîne du moteur~: aucun module du
moteur ne l'appelle. Le format xlsx est lu par l'application, hors du
package. Il ne peut appeler que \code{dictionnaire} et \code{outils}.
Ses quatre fonctions sont exportées.

\begin{center}
\small
\rowcolors{2}{tablerow}{white}
% concordance-fonctions : lecture
\begin{tabular}{@{}p{3.9cm}>{\raggedright\arraybackslash}p{3.0cm}p{8.2cm}@{}}
\rowcolor{tableheader}
\textcolor{financeblue}{\textbf{Fonction}} & \textcolor{financeblue}{\textbf{Arguments}} & \textcolor{financeblue}{\textbf{Rôle}} \\
\code{lire\_base()} & \code{dossier} & Lit une base écrite en CSV (un fichier \code{<table>.csv} par table, toutes facultatives) et la rend en liste de \code{data.frame}, dans l'ordre des tables du dictionnaire\fsp; chaque colonne connue est lue dans son type si toutes ses valeurs s'y lisent, sinon rendue entière en texte exact, sans conversion ni refus\fsp; ne s'interrompt (\code{fichier\_invalide}) que si elle ne peut pas rendre le texte exact d'un fichier\fsp; contrôle strictement la seule table \code{extraction} (section~\ref{sec:lecture-base}). \\
\code{ecrire\_base()} & \code{donnees}, \code{dossier} & Écrit une base contrôlée par \code{controler\_types()} en CSV UTF-8, un fichier par table, dans un dossier neuf ou vide, au format que relit \code{lire\_base()}\fsp; refuse (\code{liste\_invalide}) avant toute écriture la chaîne vide, hors domaine, et ce que le CSV ne relirait pas à l'identique (retour chariot, réel non fini, date non entière ou hors de quatre chiffres d'année, chaîne d'encodage inconnu ou non UTF-8)\fsp; refuse (\code{fichier\_invalide}) un dossier non vide\fsp; rend le dossier. \\
\code{ecrire\_jeu()} & \code{jeu}, \code{dossier} & Écrit un jeu simulé dans un dossier neuf ou vide, selon la disposition d'un jeu versionné~: \code{base/} (une table par fichier, au format et aux octets d'\code{ecrire\_base()}), \code{loi\_vraie.txt} et \code{manifeste.txt} (par \code{ecrire\_liste()}), \code{verite/} (une table de vérité par fichier CSV), \code{declarations.txt} si le jeu en porte, puis \code{empreintes.md5} au format de \code{md5sum -c}\fsp; refuse (\code{jeu\_invalide}) avant toute écriture, création du dossier comprise, un objet qui n'a pas la forme d'un jeu simulé ou que \code{lire\_jeu()} ne relirait pas identique, et applique avant elle aussi les refus d'\code{ecrire\_base()} et d'\code{ecrire\_liste()} (section~\ref{sec:lecture-jeu})\fsp; rend le dossier. \\
\code{lire\_jeu()} & \code{dossier} & Relit un jeu écrit par \code{ecrire\_jeu()} et le rend identique\fsp; contrôle d'abord la disposition du dossier (\code{fichier\_invalide}) et chaque empreinte de \code{empreintes.md5} (\code{empreinte\_discordante}), puis lit la base par \code{lire\_base()} (tables et nombres de lignes égaux à la volumétrie du manifeste, sinon \code{fichier\_invalide}) et les tables de vérité dans les types que décrit le manifeste. \\
\end{tabular}
\end{center}

% =============================================================================
\section{Lecture et écriture d'une base}
\label{sec:lecture-base}
% =============================================================================

\textbf{Format.} Celui du dictionnaire de la base~: un fichier
\code{<table>.csv} par table, en UTF-8 sans marque d'ordre, séparateur
virgule, point décimal, guillemets doubles, dates \code{AAAA-MM-JJ},
logiques \code{TRUE} et \code{FALSE}\fsp; le champ vide est le seul
marqueur de valeur manquante.

\textbf{Lecture «~fidélité ou refus~».} \code{lire\_base(dossier)} lit
chaque fichier en texte, sans conversion, rétablit à sa place la chaîne vide
entre guillemets, puis lit chaque colonne connue du dictionnaire dans son
type par la grammaire stricte de \code{lire\_texte\_type()}, si toutes ses
valeurs s'y lisent\fsp; sinon la colonne est rendue entière en texte exact.
La lecture typée directe de R n'est pas employée parce qu'elle convertit en
silence~: sous R~4.3.3, elle lit \code{0x10} comme le nombre~16,
\code{26-01-02} comme une date de l'an~26 et \code{T} comme la valeur
logique vraie. Une colonne entièrement vide prend ainsi son type déclaré.

\textbf{Interruptions.} \code{lire\_base()} ne s'interrompt, par le motif
\code{fichier\_invalide}, que si elle ne peut pas rendre le texte exact
d'un fichier~: dossier introuvable ou qui n'est pas un dossier\fsp; tout
fichier ou dossier autre qu'un \code{<table>.csv} d'une table du
dictionnaire\fsp; fichier mal formé (vide, marque d'ordre, UTF-8 invalide,
octet nul, dernière fin de ligne absente, retour chariot isolé hors
guillemets, guillemet au milieu d'un champ ou non fermé, nombre de champs
variable, première ligne vide, en-tête illisible, nom de colonne vide ou en
double)\fsp; retour chariot dans un champ entre guillemets, que la lecture
de R relirait comme un saut de ligne, avec la table, la colonne et les
lignes en cause. Pour la seule table \code{extraction}, qui donne lieu à un
refus d'appel en toute hypothèse, la version du dictionnaire est contrôlée
avant les autres colonnes (\code{extraction\_invalide}), puis la table
entière par \code{controler\_types()}.

\textbf{Ce qui est rendu sans refus.} Hors de la table \code{extraction},
ce qui n'a pas la forme du dictionnaire est rendu tel quel~: colonne
inconnue (en texte), colonne absente, colonne dont une valeur au moins ne
se lit pas dans son type (en texte exact), chaîne vide entre guillemets. Le
résultat d'un appel isolé n'est donc pas forcément conforme~:
\code{controler\_types()} le refuse, et le contrôle de la base entière
l'enregistre au bilan (section~\ref{sec:donnees-bilan}). Tolérances
admises, par jugement de mise en œuvre (les guillemets délimitent le champ,
ils ne portent pas le type\fsp; aucune tolérance ne change la valeur lue
d'une colonne de type caractère, et les 47~colonnes d'identifiant
\code{id\_*} du dictionnaire sont toutes de ce type, de sorte que des zéros
de tête ne peuvent confondre deux clés)~: toute valeur entre
guillemets ou non, quel que soit son type\fsp; zéros de tête d'un entier ou
d'un réel\fsp; chaîne sans guillemets, espaces conservées, \code{NA} nu lu
comme la chaîne \code{"NA"}\fsp; fins de ligne CRLF hors des champs entre
guillemets. Les colonnes sont rendues dans l'ordre du fichier, sous le
texte exact de l'en-tête, avec des noms de lignes automatiques.

\textbf{Écriture.} \code{ecrire\_base(donnees, dossier)} contrôle la base
par \code{controler\_types()}, puis écrit l'en-tête des noms de colonnes,
les chaînes entre guillemets doubles (guillemet intérieur doublé), le
manquant en champ vide, les entiers en décimal, les réels au plus court de
15, 16 ou 17 chiffres significatifs qui se relit à l'identique (17
suffisent toujours, section~\ref{sec:ecriture}), les dates
\code{AAAA-MM-JJ}, les logiques \code{TRUE} et \code{FALSE}, et des fins
de ligne \code{\textbackslash n}. Le dossier est créé s'il n'existe pas
(son parent doit exister)\fsp; s'il existe, il doit être vide, faute de quoi
un fichier ancien serait relu avec la base. La chaîne vide est refusée à
l'écriture (\code{liste\_invalide}) alors qu'elle est rendue à la
lecture~: une base écrite est conforme. \code{lire\_base(ecrire\_base(x,
d))} rend \code{x} à l'identique quand \code{x} liste ses tables dans
l'ordre du dictionnaire et porte des noms de lignes automatiques~; c'est un
constat de la batterie unitaire sur les bases testées.

% =============================================================================
\section{Écriture et relecture d'un jeu simulé}
\label{sec:lecture-jeu}
% =============================================================================

La fonction \code{ecrire\_jeu(jeu, dossier)} écrit un jeu simulé (objet
de la classe \code{jeu\_simule} que rend \code{simuler\_base()},
section~\ref{sec:simulation-jeu}) dans un dossier neuf ou vide, selon la
disposition d'un jeu versionné~:

\begin{itemize}[itemsep=2pt]
\item \code{base/}~: une table de la base par fichier \code{<table>.csv},
  aux octets qu'écrirait \code{ecrire\_base()}, dont \code{ecrire\_jeu()}
  emploie le code interne\fsp;
\item \code{verite/}~: une table de vérité par fichier \code{<nom>.csv}, au
  même format CSV, les types des colonnes étant portés par le manifeste\fsp;
\item \code{loi\_vraie.txt} et \code{manifeste.txt}, écrits par
  \code{ecrire\_liste()} (section~\ref{sec:format-liste})\fsp;
\item \code{declarations.txt}, facultatif, si le jeu porte des
  déclarations\fsp;
\item \code{empreintes.md5}~: une ligne «~\code{<md5>}, deux espaces,
  \code{<chemin relatif>}~» par fichier, triés indépendamment de la locale,
  au format que vérifie la commande \code{md5sum -c} lancée depuis le
  dossier du jeu.
\end{itemize}

Avant toute écriture, création du dossier comprise, \code{ecrire\_jeu()}
refuse (\code{jeu\_invalide}, par un message qui nomme l'élément en cause,
table et colonne comprises) un objet qui n'a pas la forme d'un jeu simulé ou
que \code{lire\_jeu()} ne relirait pas identique~:

\begin{itemize}[itemsep=2pt]
\item classe autre que \code{jeu\_simule} seule, ou attribut du jeu autre
  que ses noms et sa classe\fsp;
\item champs autres que \code{base}, \code{loi\_vraie}, \code{verite},
  \code{manifeste} dans cet ordre, puis \code{declarations} facultatif, ou
  champ \code{declarations} présent mais \code{NULL} (il ne s'écrit pas, et
  le jeu relu en serait privé)\fsp;
\item liste \code{base} ou \code{verite} qui porte un attribut autre que
  ses noms\fsp;
\item tables de vérité qui ne forment pas une liste nommée de
  \code{data.frame}, ou dont les colonnes et les types ne sont pas
  exactement ceux que décrit le manifeste\fsp;
\item tables de la base hors de l'ordre du dictionnaire (\code{lire\_base()}
  les rend dans cet ordre), ou volumétrie du manifeste différente des
  nombres de lignes de la base\fsp;
\item table de la base ou de vérité de classe autre que \code{data.frame}
  seule, dont les noms de lignes ne sont pas les noms automatiques
  \code{1..n}, qui porte un attribut de table autre que
  ses noms, sa classe et ses noms de lignes, ou dont une colonne porte un
  attribut, hors de la classe \code{Date} d'une colonne de dates~: le CSV
  ne les conserve pas.
\end{itemize}

\noindent Les refus d'\code{ecrire\_base()} s'appliquent à la base, et ses
refus de valeurs (ce que le CSV ne relirait pas à l'identique) aux tables
de vérité, au nom d'\code{ecrire\_jeu()} et de \code{verite/<nom>}\fsp; ceux
d'\code{ecrire\_liste()} (\code{liste\_invalide}, qui refuse notamment, sur
une liste, tout attribut autre que les noms) s'appliquent à la loi vraie,
au manifeste et aux déclarations. Tous sont levés avant la création du
dossier.

\code{lire\_jeu(dossier)} relit un jeu écrit par \code{ecrire\_jeu()}. Elle
contrôle d'abord la disposition (les fichiers et dossiers attendus, rien
d'autre) et la liste d'empreintes (une ligne par fichier, chacun présent une
fois), par le motif \code{fichier\_invalide}, puis l'empreinte de chaque
fichier (\code{empreinte\_discordante} au premier qui diffère). Elle lit
ensuite la base par \code{lire\_base()}, dont les tables et les nombres de
lignes doivent être ceux de la volumétrie du manifeste, les listes par
\code{lire\_liste()}, et les tables de vérité dans les types que décrit le
manifeste~; tout écart est refusé par \code{fichier\_invalide}. Un
fichier caché est refusé comme tout fichier non attendu, à la racine du jeu,
dans \code{base/} (contrôle des empreintes, puis \code{lire\_base()}) et
dans \code{verite/}. Elle
n'évalue aucun code. Que \code{lire\_jeu(ecrire\_jeu(j, d))} rende un objet
identique à \code{j} est un constat des batteries, non une démonstration.


\partie{Partie VI}{Module donnees}

% =============================================================================
\section{Rôle et fonctions du module donnees}
\label{sec:donnees}
% =============================================================================

Le module \code{donnees} contrôle la base~: détection des défauts de données
et contrôles de structure, traitements déclarés (refus, exclusion,
hypothèse par défaut, signalement), vue à une date, états des sinistres et
des dossiers à une date, segments, enjeu des lignes écartées. Il ne corrige
rien en silence. Dans l'état du code, il contrôle la base entière et la vue
à l'arrêté et chiffre l'enjeu (section~\ref{sec:donnees-bilan})\fsp; les
états des sinistres et des dossiers à une date et les segments ne sont pas
encore écrits. Il ne peut appeler que \code{dictionnaire} et
\code{outils}. Sa fonction \code{controler\_base()} est exportée.

\begin{center}
\small
\rowcolors{2}{tablerow}{white}
% concordance-fonctions : donnees
\begin{tabular}{@{}p{3.9cm}>{\raggedright\arraybackslash}p{3.0cm}p{8.2cm}@{}}
\rowcolor{tableheader}
\textcolor{financeblue}{\textbf{Fonction}} & \textcolor{financeblue}{\textbf{Arguments}} & \textcolor{financeblue}{\textbf{Rôle}} \\
\code{controler\_base()} & \code{donnees}, \code{arrete}, \code{declarations = NULL}, \code{conventions = NULL} & Contrôle la base entière et la vue à la date d'arrêté, sans rien modifier, sous les déclarations du sous-schéma \code{donnees} du plan et la section \code{conventions} du plan, et rend le bilan en onze éléments (section~\ref{sec:donnees-bilan})\fsp; chaque contrôle prévu et reporté figure dans les contrôles non exécutés avec la cause \code{non\_implemente}\fsp; refuse l'appel pour un arrêté absent ou mal formé, une liste mal formée, une table inconnue ou obligatoire absente, une anomalie de la table \code{extraction}, un arrêté postérieur à l'extraction, une déclaration ou une convention mal formée ou hors de sa liste fermée. \\
\end{tabular}
\end{center}

% =============================================================================
\section{Bilan de la base}
\label{sec:donnees-bilan}
% =============================================================================

\code{controler\_base(donnees, arrete, declarations, conventions)}
applique la spécification du module \code{donnees} visée par le mainteneur
(\code{docs/specifications/module-donnees-v1.md}, version~1.1), sous le
dictionnaire de la base en version~1.5. La date d'arrêté est obligatoire,
sans valeur par défaut. \code{declarations} est le sous-schéma
\code{donnees} du plan (rubriques \code{exclusions}, \code{hypotheses},
\code{seuils}, \code{evaluations\_gestion})\fsp; \code{conventions} est la
section \code{conventions} du plan, dont seules les entrées \code{age} et
\code{limite\_age} sont lues ici. L'une et l'autre peuvent être omises
(\code{NULL})~: cela se lit «~aucune déclaration fournie~», et non comme
une valeur par défaut.

\textbf{Ordre d'exécution.} (1)~Interruptions (refus d'appel)\fsp; (2)~forme,
(3)~domaine, (4)~clés et refus induits, qui forment la \textbf{règle~B}, sur
la base entière et indépendante de l'arrêté\fsp; (5)~détection des défauts
D02b et D03g sur la base entière (\textbf{règle~P}, préalables de la
vue)\fsp; (6)~construction de la vue à l'arrêté\fsp; (7)~détection, sur la
vue, des autres défauts traités par exclusion, hypothèse par défaut ou
signalement (\textbf{règle~V})\fsp; (8)~enjeu des lignes écartées. Chaque
étape suppose les précédentes~: une clé mal typée ne se joint pas, et la vue
suppose des clés résolues. Les règles~P et~V et la vue ne lisent ni une
table ni une colonne refusées en règle~B.

\textbf{Traitements.} Chaque anomalie est enregistrée avec son traitement~:
refus~(R) d'une table ou d'une colonne, exclusion~(X), hypothèse par
défaut~(H) ou signalement~(S)\fsp; aucune donnée n'est modifiée. Un code
de contrôle de structure de portée \code{ligne} est traité en
signalement, tout autre en refus (règle provisoire, fixée par l'ADR~0010).
Un constat qui admet un choix est placé selon les déclarations, dans un
seul élément par famille de calcul ou par référentiel\fsp; sans déclaration,
il va dans \code{declarations\_requises}.

\textbf{Le bilan.} Il est rendu en liste nommée de onze éléments, dans cet
ordre~:

\begin{center}
\small
\rowcolors{2}{tablerow}{white}
\begin{tabular}{@{}p{4.2cm}p{10.9cm}@{}}
\rowcolor{tableheader}
\textcolor{financeblue}{\textbf{Élément}} & \textcolor{financeblue}{\textbf{Contenu}} \\
\code{controles} & Toute détection~: code, table, colonne, détection (exacte ou présomptive), nombre de lignes, identifiants. \\
\code{controles\_non\_executes} & Contrôle empêché, une ligne par (code, table, colonne), avec sa cause (ci-dessous). \\
\code{refus\_tables} & Refus~R~: table, colonne, motif, portée effective (\code{colonne} ou \code{table}), origine d'un refus induit, nombre de lignes, lignes visibles à l'arrêté (\code{NA} si indéterminé), identifiants, dossiers rattachés visibles, postes arrêtés (\code{NA} dans l'état du code). \\
\code{exclusions} & Exclusions déclarées~: motif, famille de calcul, action (\code{exclure} ou \code{conserver}), nombre de lignes et de dossiers. \\
\code{hypotheses\_defaut} & Hypothèses par défaut déclarées~: motif, référentiel, poste, valeur, nombre. \\
\code{signalements} & Constats traités en signalement, ou dont le signalement est déclaré~: mêmes colonnes que \code{controles}. \\
\code{declarations\_requises} & Déclaration attendue par un constat ou par un contrôle~: rubrique (\code{exclusions}, \code{hypotheses}, \code{seuils} ou \code{conventions}), motif, clé (famille, référentiel, seuil ou entrée de convention), traitements ou valeurs admis, nombre de lignes. \\
\code{enjeu} & Enjeu des lignes refusées ou exclues, en table longue~: une ligne par table, colonne, motif, famille, action, garantie, poste et composante, avec son unité, son montant, ses nombres de dossiers et de valeurs, et la cause d'une composante empêchée. \\
\code{vue} & Dates d'arrêté et d'extraction, lignes de visibilité non déterminée, et indicateur \code{provisoire}~: vrai s'il y en a, ou si une règle~P n'a pas pu s'exécuter sur une table que lit la vue (les comptes de la règle~V sont alors provisoires). \\
\code{segments} & Vide dans l'état du code (les segmentations se déclarent avec le plan). \\
\code{tables\_absentes} & Tables facultatives absentes de la base. \\
\end{tabular}
\end{center}

Chaque \code{data.frame} du bilan est trié dans l'ordre des tables du
dictionnaire, puis des colonnes (\code{NA} en tête), puis du code\fsp; les
identifiants sont les dix premières valeurs de la clé primaire, triées
indépendamment de la locale. Les composantes de l'enjeu forment une liste
fermée, \code{COMPOSANTES\_ENJEU}~: huit composantes, dix couples
(composante, unité), aucune agrégation ne mêlant deux couples.

\textbf{Contrôles non exécutés et causes.} La colonne \code{cause} de
\code{controles\_non\_executes} prend l'une des valeurs suivantes~:

\begin{itemize}[itemsep=2pt]
\item le motif d'un refus d'une colonne lue (\textbf{cause directe})~:
  code de contrôle (\code{controles.csv}), code de défaut
  (\code{defauts.csv}) ou sous-cas lettré du défaut D10 (D10a à D10f et
  D10h, la lettre~g n'étant attribuée à aucun sous-cas), liste fermée\fsp;
\item \code{refus\_induit}~: table refusée parce qu'une de ses clés
  étrangères vise une table refusée, ou, à un arrêté antérieur à
  l'extraction, parce qu'une règle~P n'a pas pu s'exécuter sur une table
  que lit la vue\fsp;
\item \code{non\_implemente}~: contrôle prévu par la spécification et
  reporté (ci-dessous)\fsp;
\item \code{declaration\_absente}~: contrôle qui attend une convention ou un
  seuil non déclarés\fsp;
\item \code{non\_lisible}~: sous-cas qu'aucune colonne de la base ne permet
  de lire\fsp;
\item \code{age\_non\_defini}~: date de référence antérieure à la
  naissance\fsp;
\item \code{parametre\_non\_defini}~: limite d'âge non déterminable pour la
  fin de charge d'un enfant.
\end{itemize}

Une composante de l'enjeu empêchée par un refus y figure aussi, sous le
code de la composante, avec la table et la colonne lues. Une composante
dont la visibilité à l'arrêté n'est pas déterminée n'y figure pas~: sa
cause, \code{visibilite\_non\_determinee}, se lit dans l'élément
\code{enjeu}. Une clé (code, table, colonne) présente à la fois dans
\code{controles} et dans \code{controles\_non\_executes} est une détection
partielle, dont le nombre est un minorant.

\textbf{Priorité des causes.} Quand plusieurs causes visent une même clé
(code, table, colonne), parce que plusieurs règles, emplois ou colonnes
lues l'empêchent, \code{controles\_non\_executes} n'en garde qu'une ligne,
qui porte la cause de plus petit rang (décisions du mainteneur du
8~novembre~2026, M185, et du 9~novembre~2026, M194 et M195)~:

\begin{enumerate}[itemsep=2pt]
\item \code{non\_implemente}\fsp;
\item \code{non\_lisible}\fsp;
\item cause directe, quelle qu'elle soit\fsp;
\item \code{age\_non\_defini}\fsp;
\item \code{parametre\_non\_defini}\fsp;
\item \code{declaration\_absente}\fsp;
\item \code{refus\_induit}.
\end{enumerate}

À rang égal (deux causes directes distinctes), la première
émise l'emporte. Un refus induit n'est donc retenu que si la clé n'a
aucune autre cause, et la priorité de la cause directe sur tout refus
induit, règle~B comprise (M185), est conservée. L'ordre est un jugement
d'expert, motivé ainsi~: (a)~la cause retenue ne dépend pas de l'ordre
d'exécution du code, de sorte que réordonner deux règles ou deux emplois
ne change pas le bilan (l'alternative écartée, la première cause émise,
faisait dépendre le bilan de cet agencement)\fsp; (b)~passent d'abord les
causes que rien ne lève côté utilisateur (l'outil, puis le
dictionnaire), puis le défaut de l'extraction, puis la grandeur que la
base ne détermine pas\fsp; (c)~les deux derniers rangs vont aux causes
dont l'information reste au bilan quand elles ne sont pas retenues~: une
déclaration absente figure toujours dans \code{declarations\_requises},
le refus d'origine d'un refus induit dans \code{refus\_tables}, alors
qu'une cause directe, un âge non défini ou un paramètre non défini
masqués seraient perdus.

Dans le code, la fonction interne \code{.dedoublonner\_non\_executes()}
applique cet ordre, lu dans deux constantes internes~:
\code{.RANGS\_CAUSES}, qui donne les rangs des six causes nommées
ci-dessus, et \code{.RANG\_CAUSE\_DIRECTE}, qui vaut~3. Le rang d'une
cause est cherché dans \code{.RANGS\_CAUSES} avant de tester si elle est
directe, parce que \code{refus\_induit} est aussi un code de
\code{controles.csv}. Une cause qui n'est ni rangée ni dans la liste
fermée des causes directes arrête le calcul, avec un message qui
commence par «~erreur interne de provie~»~: elle n'est jamais rangée en
silence, et toute cause nouvelle doit recevoir un rang dans
\code{.RANGS\_CAUSES}. \code{visibilite\_non\_determinee} n'a pas de
rang~: elle n'est jamais une cause de \code{controles\_non\_executes}, et
la fonction s'arrête si on la lui présente. La batterie unitaire vérifie
l'ordre sur chaque paire de rangs, dans les deux ordres d'émission, le
départage à rang égal, l'arrêt sur une cause inconnue (dont une lettre
ajoutée à un code déjà lettré, comme D01ab, et un sous-cas non déclaré,
comme D05z), et, sur l'ensemble des appels qu'elle fait, que toute cause
émise est rangée ou directe, que les cinq causes non directes autres que
\code{refus\_induit} sont émises, et que
\code{visibilite\_non\_determinee} ne l'est jamais.

\textbf{Âge invraisemblable et fin de charge.} Le défaut D02c, pour un
âge, se détecte sur la vue, pour cinq emplois (assuré au sinistre, assuré
à l'affiliation, membre de la famille, désigné, ayant droit), sous la
convention \code{conventions\$age} et les bornes déclarées
\code{seuils\$d02c\_age\_min} et \code{seuils\$d02c\_age\_max} de chaque
emploi\fsp; traitement signalement et exclusion. Le défaut D07, sous-cas~(c)
(fin de charge non enregistrée), vise une rente d'éducation en service ou
un dossier de double effet ni clos ni de second décès retenu, dont chaque
enfant à charge a atteint sa limite d'âge à l'arrêté. La limite se lit
sous deux entrées déclarées~: la règle, \code{atteinte} ou
\code{depassement} (entrée \code{conventions\$limite\_age\$regle}), et la
convention d'âge, \code{revolu} ou \code{millesime} (entrée
\code{conventions\$limite\_age\$age})\fsp; traitement hypothèse par défaut ou
signalement. Sans la convention ou les bornes, le contrôle n'est pas
exécuté (\code{declaration\_absente}) et la déclaration manquante figure
dans \code{declarations\_requises}.

\textbf{État partiel déclaré.} Le bilan dit lui-même qu'il est partiel~:
les contrôles prévus par la spécification et reportés à l'issue~34 du
dépôt par décision du mainteneur (M181) figurent, pour chaque table
présente, dans les contrôles non exécutés, avec la cause
\code{non\_implemente}. Ils sont énumérés en quatorze lignes (code, table,
colonne) sur dix codes par la constante \code{CONTROLES\_NON\_IMPLEMENTES}~: D03d\fsp; D07 sur \code{evenements} (sous-cas~(a)
et~(b))\fsp; D08\fsp; D09 aux seuils, sur cinq tables\fsp; D10d\fsp; D10e\fsp;
D10f\fsp; D11\fsp; D12\fsp; D14. Aucun seuil de ces contrôles n'est
attendu dans \code{declarations\_requises}.

Les constantes du module sont citées d'après le code (\code{R/donnees.R})~:

\begin{center}
\small
\rowcolors{2}{tablerow}{white}
\begin{tabular}{@{}p{5.6cm}p{9.5cm}@{}}
\rowcolor{tableheader}
\textcolor{financeblue}{\textbf{Constante}} & \textcolor{financeblue}{\textbf{Valeur et objet}} \\
\code{FAMILLES\_CALCUL} & \code{engagement}, \code{triangle}, \code{risque\_en\_cours}, \code{regle\_comptable}~: familles d'une exclusion déclarée. \\
\code{REFERENTIELS} & \code{comptes\_sociaux}, \code{solvabilite\_2}~: référentiels d'une hypothèse par défaut. \\
\code{ACTIONS\_EXCLUSION} & \code{exclure}, \code{conserver}~: actions d'une exclusion déclarée. \\
\code{RUBRIQUES\_DECLARATIONS} & \code{exclusions}, \code{hypotheses}, \code{seuils}, \code{evaluations\_gestion}~: rubriques admises de \code{declarations}. \\
\code{ENTREES\_CONVENTIONS} & \code{age}, \code{instant\_date}, \code{repartition\_annee}, \code{fractionnement}, \code{limite\_age}~: entrées admises de \code{conventions}\fsp; seules \code{age} et \code{limite\_age} sont contrôlées et lues. \\
\code{REGLES\_LIMITE\_AGE} & \code{atteinte}, \code{depassement}~: règle de limite d'âge du défaut D07~(c). \\
\code{CONVENTIONS\_LIMITE\_AGE} & \code{revolu}, \code{millesime}~: conventions d'âge admises pour le défaut D07~(c). \\
\code{COMPOSANTES\_ENJEU} & Tableau des composantes de l'enjeu (composante, unité, définition), dix lignes. \\
\code{CAUSE\_NON\_IMPLEMENTE} & \code{"non\_implemente"}~: cause d'un contrôle reporté. \\
\code{CAUSE\_NON\_LISIBLE} & \code{"non\_lisible"}~: cause d'un sous-cas illisible dans la base. \\
\code{CONTROLES\_NON\_IMPLEMENTES} & Tableau des contrôles reportés (code, table, colonne), quatorze lignes. \\
\end{tabular}
\end{center}


\partie{Partie VII}{Module tables}

% =============================================================================
\section{Rôle et fonctions du module tables}
\label{sec:tables}
% =============================================================================

Le module \code{tables} traite les tables en entrée~: formats admis,
conversion vers une représentation commune, catalogue des tables
embarquées, courbes des taux sans risque de l'EIOPA, série de taux des
plafonds réglementaires, plafonds réglementaires. Le calcul de l'âge n'en
fait pas partie~: il est dans le module \code{outils}, qui sert aussi
\code{donnees} et le simulateur (section~\ref{sec:age}). Le module \code{tables} ne
peut appeler que \code{outils}. Ce module n'a pas encore de fonction.

\begin{center}
\small
\rowcolors{2}{tablerow}{white}
% concordance-fonctions : tables
\begin{tabular}{@{}p{3.9cm}>{\raggedright\arraybackslash}p{3.0cm}p{8.2cm}@{}}
\rowcolor{tableheader}
\textcolor{financeblue}{\textbf{Fonction}} & \textcolor{financeblue}{\textbf{Arguments}} & \textcolor{financeblue}{\textbf{Rôle}} \\
\multicolumn{3}{@{}l@{}}{Aucune fonction.} \\
\end{tabular}
\end{center}


\partie{Partie VIII}{Module garanties}

% =============================================================================
\section{Rôle et fonctions du module garanties}
\label{sec:garanties}
% =============================================================================

Le module \code{garanties} lit et valide le catalogue des configurations de
garanties~: états, transitions, flux, postes, usages de tables, liens entre
garanties. Une configuration est une donnée, jamais une fonction. Le module
ne peut appeler que \code{dictionnaire} et \code{outils}. Ce module n'a pas
encore de fonction.

\begin{center}
\small
\rowcolors{2}{tablerow}{white}
% concordance-fonctions : garanties
\begin{tabular}{@{}p{3.9cm}>{\raggedright\arraybackslash}p{3.0cm}p{8.2cm}@{}}
\rowcolor{tableheader}
\textcolor{financeblue}{\textbf{Fonction}} & \textcolor{financeblue}{\textbf{Arguments}} & \textcolor{financeblue}{\textbf{Rôle}} \\
\multicolumn{3}{@{}l@{}}{Aucune fonction.} \\
\end{tabular}
\end{center}


\partie{Partie IX}{Module engagement}

% =============================================================================
\section{Rôle et fonctions du module engagement}
\label{sec:engagement}
% =============================================================================

Le module \code{engagement} porte la famille de calcul «~engagement~»~:
échéanciers des flux futurs d'un engagement connu, dossier par dossier, et
leur valeur actuelle probable sous le modèle multi-états. Il ne peut
appeler que \code{tables}, \code{garanties}, \code{donnees} et
\code{outils}. Ce module n'a pas encore de fonction.

\begin{center}
\small
\rowcolors{2}{tablerow}{white}
% concordance-fonctions : engagement
\begin{tabular}{@{}p{3.9cm}>{\raggedright\arraybackslash}p{3.0cm}p{8.2cm}@{}}
\rowcolor{tableheader}
\textcolor{financeblue}{\textbf{Fonction}} & \textcolor{financeblue}{\textbf{Arguments}} & \textcolor{financeblue}{\textbf{Rôle}} \\
\multicolumn{3}{@{}l@{}}{Aucune fonction.} \\
\end{tabular}
\end{center}


\partie{Partie X}{Module triangle}

% =============================================================================
\section{Rôle et fonctions du module triangle}
\label{sec:triangle}
% =============================================================================

Le module \code{triangle} porte la famille de calcul «~triangle~»~:
reconstitution des triangles à partir de la base sinistres, estimateurs,
diagnostics. Il ne peut appeler que \code{garanties}, \code{donnees} et
\code{outils}. Ce module n'a pas encore de fonction.

\begin{center}
\small
\rowcolors{2}{tablerow}{white}
% concordance-fonctions : triangle
\begin{tabular}{@{}p{3.9cm}>{\raggedright\arraybackslash}p{3.0cm}p{8.2cm}@{}}
\rowcolor{tableheader}
\textcolor{financeblue}{\textbf{Fonction}} & \textcolor{financeblue}{\textbf{Arguments}} & \textcolor{financeblue}{\textbf{Rôle}} \\
\multicolumn{3}{@{}l@{}}{Aucune fonction.} \\
\end{tabular}
\end{center}


\partie{Partie XI}{Module risque\_en\_cours}

% =============================================================================
\section{Rôle et fonctions du module risque\_en\_cours}
\label{sec:risque-en-cours}
% =============================================================================

Le module \code{risque\_en\_cours} porte la famille de calcul «~risque en
cours~»~: limites du contrat et provision pour risques croissants (PRC),
hors de la première version. Il ne peut appeler que \code{engagement},
\code{tables}, \code{garanties}, \code{donnees} et \code{outils}. Ce module
n'a pas encore de fonction.

\begin{center}
\small
\rowcolors{2}{tablerow}{white}
% concordance-fonctions : risque_en_cours
\begin{tabular}{@{}p{3.9cm}>{\raggedright\arraybackslash}p{3.0cm}p{8.2cm}@{}}
\rowcolor{tableheader}
\textcolor{financeblue}{\textbf{Fonction}} & \textcolor{financeblue}{\textbf{Arguments}} & \textcolor{financeblue}{\textbf{Rôle}} \\
\multicolumn{3}{@{}l@{}}{Aucune fonction.} \\
\end{tabular}
\end{center}


\partie{Partie XII}{Module regles\_comptables}

% =============================================================================
\section{Rôle et fonctions du module regles\_comptables}
\label{sec:regles-comptables}
% =============================================================================

Le module \code{regles\_comptables} porte la catégorie «~règle
comptable~»~: postes calculés par application d'une règle déclarée du plan
(dotation, reprise, prorata temporis, ratio, seuil) à des grandeurs
reconstituées de la base, tels la provision pour égalisation, la provision
pour primes non acquises (PPNA) et la provision pour risques en cours
(PREC). Une règle comptable n'est pas une famille de calcul. Le module ne
peut appeler que \code{garanties}, \code{tables}, \code{donnees} et
\code{outils}. Ce module n'a pas encore de fonction.

\begin{center}
\small
\rowcolors{2}{tablerow}{white}
% concordance-fonctions : regles_comptables
\begin{tabular}{@{}p{3.9cm}>{\raggedright\arraybackslash}p{3.0cm}p{8.2cm}@{}}
\rowcolor{tableheader}
\textcolor{financeblue}{\textbf{Fonction}} & \textcolor{financeblue}{\textbf{Arguments}} & \textcolor{financeblue}{\textbf{Rôle}} \\
\multicolumn{3}{@{}l@{}}{Aucune fonction.} \\
\end{tabular}
\end{center}


\partie{Partie XIII}{Module projection}

% =============================================================================
\section{Rôle et fonctions du module projection}
\label{sec:projection}
% =============================================================================

Le module \code{projection} portera, dans la deuxième version, la famille
de calcul «~projection~»~: projection actif-passif sous scénarios
économiques. Il ne peut appeler que \code{engagement}, \code{tables},
\code{garanties}, \code{donnees} et \code{outils}. Ce module n'a pas encore
de fonction.

\begin{center}
\small
\rowcolors{2}{tablerow}{white}
% concordance-fonctions : projection
\begin{tabular}{@{}p{3.9cm}>{\raggedright\arraybackslash}p{3.0cm}p{8.2cm}@{}}
\rowcolor{tableheader}
\textcolor{financeblue}{\textbf{Fonction}} & \textcolor{financeblue}{\textbf{Arguments}} & \textcolor{financeblue}{\textbf{Rôle}} \\
\multicolumn{3}{@{}l@{}}{Aucune fonction.} \\
\end{tabular}
\end{center}


\partie{Partie XIV}{Module plan}

% =============================================================================
\section{Rôle et fonctions du module plan}
\label{sec:plan}
% =============================================================================

Le module \code{plan} valide et met sous forme canonique le plan de calcul,
établit la fiche de calcul et sérialise l'un et l'autre en texte. Il ne
peut appeler que \code{tables}, \code{garanties}, \code{dictionnaire},
\code{donnees} et \code{outils}. Ce module n'a pas encore de fonction.

\begin{center}
\small
\rowcolors{2}{tablerow}{white}
% concordance-fonctions : plan
\begin{tabular}{@{}p{3.9cm}>{\raggedright\arraybackslash}p{3.0cm}p{8.2cm}@{}}
\rowcolor{tableheader}
\textcolor{financeblue}{\textbf{Fonction}} & \textcolor{financeblue}{\textbf{Arguments}} & \textcolor{financeblue}{\textbf{Rôle}} \\
\multicolumn{3}{@{}l@{}}{Aucune fonction.} \\
\end{tabular}
\end{center}


\partie{Partie XV}{Module rapport}

% =============================================================================
\section{Rôle et fonctions du module rapport}
\label{sec:rapport}
% =============================================================================

Le module \code{rapport} établit la représentation canonique du résultat et
ses empreintes, les quantités restituées, le tableau de passage, la
comparaison de deux résultats et le rapport figé. Il peut appeler tous les
modules qui le précèdent dans le tableau de la section~\ref{sec:modules},
sauf \code{lecture}. Ce module n'a pas encore de fonction.

\begin{center}
\small
\rowcolors{2}{tablerow}{white}
% concordance-fonctions : rapport
\begin{tabular}{@{}p{3.9cm}>{\raggedright\arraybackslash}p{3.0cm}p{8.2cm}@{}}
\rowcolor{tableheader}
\textcolor{financeblue}{\textbf{Fonction}} & \textcolor{financeblue}{\textbf{Arguments}} & \textcolor{financeblue}{\textbf{Rôle}} \\
\multicolumn{3}{@{}l@{}}{Aucune fonction.} \\
\end{tabular}
\end{center}


\partie{Partie XVI}{Module simulation}

% =============================================================================
\section{Rôle et fonctions du module simulation}
\label{sec:simulation}
% =============================================================================

Le module \code{simulation} est le simulateur~: il tire une base brute
fictive sous une loi vraie et une graine, puis y injecte les défauts de
données paramétrés. Il rend un objet en mémoire et n'écrit aucun fichier~:
l'écriture et la relecture d'un jeu sont dans le module \code{lecture}
(section~\ref{sec:lecture-jeu}). Il est hors de la chaîne du moteur~: aucun
module du moteur ne l'appelle. Il peut appeler \code{dictionnaire},
\code{garanties}, \code{tables} et \code{outils}\fsp; dans l'état du code,
il n'appelle que \code{dictionnaire} et \code{outils}~: les tables de la loi
vraie y sont passées par valeurs, et la configuration des garanties liées
est un paramètre de la loi vraie. Sa fonction \code{simuler\_base()} et
les deux méthodes de la classe \code{jeu\_simule} sont exportées.

\begin{center}
\small
\rowcolors{2}{tablerow}{white}
% concordance-fonctions : simulation
\begin{tabular}{@{}p{3.9cm}>{\raggedright\arraybackslash}p{3.0cm}p{8.2cm}@{}}
\rowcolor{tableheader}
\textcolor{financeblue}{\textbf{Fonction}} & \textcolor{financeblue}{\textbf{Arguments}} & \textcolor{financeblue}{\textbf{Rôle}} \\
\code{simuler\_base()} & \code{loi\_vraie}, \code{graines}, \code{date\_extraction}, \code{nom}, \code{description}, \code{declarations = NULL}, \code{date\_visa = NULL} & Tire une base brute fictive sous la loi vraie, en flux séparés (population, famille, sinistres, information), chacun sous sa graine, puis injecte chaque défaut de données paramétré sous la graine des défauts augmentée de son rang\fsp; rend un objet \code{jeu\_simule} (base, loi vraie, tables de vérité dont le registre des défauts, manifeste avec la version du simulateur, déclarations facultatives) sans écrire de fichier\fsp; refuse (\code{jeu\_invalide}), avant tout tirage, une loi vraie, des graines ou des arguments mal formés, une loi vraie à clé inconnue ou à valeur hors de son domaine, et une entrée qui ne satisfait pas les deux préconditions du tirage (section~\ref{sec:simulation-tirage}). \\
\code{format.jeu\_simule()} & \code{x}, \code{...} & Décrit un jeu simulé en lignes de texte, par des décomptes seuls~: nom, date d'extraction, versions du dictionnaire, du simulateur et du package, date de visa, lignes par table, lignes du registre des défauts par code, déclarations. \\
\code{print.jeu\_simule()} & \code{x}, \code{...} & Imprime \code{format(x)} et rend le jeu, invisible. \\
\end{tabular}
\end{center}

% =============================================================================
\section{Tirage d'une base brute fictive}
\label{sec:simulation-tirage}
% =============================================================================

\textbf{Modèle.} Les assurés sont affiliés à des contrats collectifs
obligatoires qui portent trois garanties liées~: décès, incapacité,
invalidité. En activité, le décès et l'arrêt de travail sont en risques
concurrents, par pas de 30~jours, sous des taux annuels par âge. Un arrêt
suit un modèle mensuel par ancienneté (reprise, décès, passage en
invalidité)\fsp; il passe en invalidité ou s'éteint par fin de droits au
terme de la durée maximale d'indemnisation, et prend fin à l'âge de fin de
service. L'invalidité suit des taux mensuels constants (reprise, décès). Le
décès d'un assuré en activité, en arrêt ou en invalidité ouvre un sinistre décès
(maintien de la garantie décès pendant l'arrêt)\fsp; un arrêt peut être
classé sans suite. Chaque paramètre est lu dans la \textbf{loi vraie},
liste nommée de version \code{VERSION\_LOI\_VRAIE}, où il est un jugement
d'expert déclaré~: le simulateur n'a aucune valeur par défaut, et une loi
vraie à laquelle manque une feuille est refusée (\code{jeu\_invalide}). Une
probabilité annuelle de décès $q_x = 1$ est admise~: le décès est alors certain
dans le pas, où il prime sur l'arrêt de travail (limite du partage entre les
deux risques concurrents quand l'intensité de décès devient infinie).

\textbf{Contrôles d'entrée.} \code{simuler\_base()} refuse
(\code{jeu\_invalide}), avant tout tirage~: une loi vraie d'une autre
version, à clé inconnue ou mal nommée, ou dont une valeur sort de son
domaine (probabilité hors de l'intervalle $[0,1]$, taux, cadence, âge, durée ou montant
négatif, nombre de contrats ou d'assurés par contrat inférieur à~1, âge
d'entrée minimal supérieur au maximal ou maximal au moins égal à l'âge de
retraite, arrondi du salaire nul, valeur non finie), ou dont les tables
sont incohérentes (table de maintien en incapacité dont les mois ne sont
pas 0, 1, 2… ou dont les probabilités d'un mois ont une somme supérieure
à~1, ou plus courte que la durée maximale d'indemnisation\fsp; tables de
mortalité ou d'incidence d'arrêt à âges non consécutifs ou à colonnes de
longueurs différentes\fsp; probabilités des catégories d'invalidité de
somme différente de~1\fsp; intensités de défauts qui ne donnent pas
exactement une valeur de $[0,1]$ par défaut simulé)\fsp; une date
d'extraction non entière ou non finie, ou qui ne suit pas strictement le
début de l'historique. S'y ajoutent deux préconditions du tirage, contrôlées
plutôt qu'imposées par une borne, pour que les tirages des jeux admis ne
changent pas~: (i)~l'effet du dernier contrat possible (premier jour
du mois situé, après celui du début de l'historique, au nombre maximal de
mois de décalage d'effet, \code{decalage\_effet\_mois\_max}) ne suit pas la
date d'extraction, faute de quoi un contrat, ses affiliations et ses
entrées seraient datés après elle\fsp; (ii)~les uniformes d'invalidité,
en nombre fixe par arrêt, couvrent le plus long séjour possible en
invalidité, de l'âge d'entrée minimal à l'âge de retraite, ce qui garantit
tout écart d'au plus 68~ans entre ces deux âges. Pour un écart entier~$d$, la
condition $365{,}25\,d + 2 \le 360\,(d+1)$ équivaut à
$d \le 358/5{,}25 \approx 68{,}19$~: 68~ans est admis et 69~refusé (calcul
exact, vérifié par \code{simuler\_base()} sur la loi vraie du jeu
\code{v1-modeste-01} aux seuls âges modifiés). L'âge d'entrée minimal étant
réel, au-delà de 68,24~ans un écart non entier est admis ou refusé selon sa
partie fractionnaire, le nombre d'uniformes étant arrondi au mois
supérieur.

\textbf{Flux séparés.} Le jeu propre est tiré en quatre flux, dans cet
ordre~: \code{population} (contrats, assurés, affiliations),
\code{famille} (conjoints vrais), \code{sinistres} (trajectoires vraies des
assurés et des arrêts) et \code{information} (délais de déclaration, de
saisie, d'instruction et de règlement). Chaque flux est tiré par un seul
appel d'\code{avec\_graine()} (section~\ref{sec:graine}), sous sa graine,
en tableaux d'uniformes de taille fixe par unité (contrat, assuré, arrêt,
sinistre), transformés ensuite par inversion. Le nombre d'uniformes d'un
flux ne dépend que du nombre d'unités, issu des flux amont, et de bornes
fixes, jamais des paramètres du flux ni d'un flux aval~: le tirage d'une
unité ne dépend pas de ce qui arrive aux autres, et le flux
\code{information} ne change aucun tirage des flux qui le précèdent. Les
bornes de taille (ainsi le nombre maximal d'arrêts par assuré) sont des
paramètres de la loi vraie.

\textbf{Défauts injectés.} Chaque défaut du catalogue
\code{DEFAUTS\_SIMULES} d'intensité non nulle est ensuite injecté sous la
graine des défauts augmentée de son rang dans le catalogue (et non dans la
liste des défauts actifs). L'intensité, lue dans la loi vraie, est la
probabilité d'injection par ligne éligible. Les lignes éligibles et les
tirages d'un défaut ne lisent que le jeu propre, et deux défauts d'une même
table touchent des colonnes distinctes ou des lignes disjointes~: activer
ou désactiver un défaut ne change donc ni le jeu propre ni les autres
défauts. Chaque modification est consignée au registre des défauts (table
de vérité \code{defauts}, colonnes \code{COLONNES\_REGISTRE}~: code,
paramètre, rang, table, clé, colonne, valeur propre, valeur injectée). Le
catalogue compte douze défauts du dictionnaire, chacun limité à un
sous-cas~; les autres défauts ne sont pas simulés, et leur paramètre est
refusé dans la loi vraie.

\begin{center}
\small
\rowcolors{2}{tablerow}{white}
\begin{tabular}{@{}p{0.8cm}p{1.1cm}p{4.8cm}p{2.3cm}>{\raggedright\arraybackslash}p{4.7cm}@{}}
\rowcolor{tableheader}
\textcolor{financeblue}{\textbf{Rang}} & \textcolor{financeblue}{\textbf{Code}} & \textcolor{financeblue}{\textbf{Paramètre}} & \textcolor{financeblue}{\textbf{Table}} & \textcolor{financeblue}{\textbf{Lignes éligibles et effet}} \\
1 & D01b & \code{doublon\_contenu} & \code{cotisations} & Encaissements antérieurs à l'extraction~: copie sous un nouvel identifiant, datée après l'original, au plus à l'extraction. \\
2 & D02a & \code{date\_ordre} & \code{prestations} & Indemnités journalières~: début et fin de période permutés. \\
3 & D02b & \code{saisie\_anterieure} & \code{evenements} & Événements d'un type non programmable~: saisie avant la date d'effet. \\
4 & D03a & \code{manquant\_date\_naissance} & \code{assures} & Date de naissance effacée. \\
5 & D03b & \code{manquant\_sexe} & \code{assures} & Sexe effacé. \\
6 & D03c & \code{manquant\_salaire} & \code{affiliations} & Salaire de référence effacé. \\
7 & D03g & \code{manquant\_date\_saisie} & \code{prestations} & Date de saisie effacée. \\
8 & D04b & \code{orpheline} & \code{cotisations} & Émissions~: garantie souscrite remplacée par un identifiant inconnu. \\
9 & D05 & \code{clos\_sans\_motif} & \code{evenements} & Sorties sans suite~: motif effacé. \\
10 & D06 & \code{cloture\_sans\_justificatif} & \code{evenements} & Sorties pour reprise du travail ou décès~: justificatif effacé. \\
11 & D09 & \code{montant\_aberrant} & \code{prestations} & Paiements~: montant de signe inversé. \\
12 & D13 & \code{instruction\_absente} & \code{sinistres} & Sinistres décès dont le capital est réglé~: date d'instruction effacée. \\
\end{tabular}
\end{center}

% =============================================================================
\section{Jeu simulé et jeux versionnés}
\label{sec:simulation-jeu}
% =============================================================================

\textbf{Jeu simulé.} \code{simuler\_base()} rend un objet de classe
\code{jeu\_simule}, liste nommée~:

\begin{itemize}[itemsep=2pt]
\item \code{base}~: les tables obligatoires du dictionnaire, dans son
  ordre, colonnes dans l'ordre du dictionnaire, lignes saisies au plus tard
  à la date d'extraction\fsp;
\item \code{loi\_vraie}~: la loi vraie reçue, inchangée\fsp;
\item \code{verite}~: ce que l'organisme ignore, en cinq tables de vérité
  décrites par la constante \code{SCHEMA\_VERITE}, seule définition de
  leurs colonnes et de leurs types~: \code{flux\_futurs} (prestations
  réglées après l'extraction, avec la date de survenance du sinistre et,
  pour un décès survenu après l'extraction au cours d'un arrêt de travail ou d'une invalidité,
  l'identifiant de l'arrêt d'origine), \code{etats\_vrais} (trajectoire
  vraie de chaque arrêt), \code{familles\_vraies} (conjoints vrais),
  \code{sinistres\_non\_declares} (sinistres survenus au plus tard à
  l'extraction et saisis après) et \code{defauts} (registre des
  défauts)\fsp;
\item \code{manifeste}~: dans cet ordre, nom, version du dictionnaire,
  version du package qui a produit le jeu, version du simulateur, graines,
  date d'extraction, volumétrie par table, description, date de visa, et
  colonnes des tables de vérité\fsp;
\item \code{declarations}, facultatif~: les déclarations et conventions de
  \code{controler\_base()} qui conviennent au jeu, portées sans être
  lues.
\end{itemize}

Les graines forment une liste nommée de cinq entiers, une par flux et une
pour les défauts (\code{GRAINES\_SIMULATION}). À loi vraie, graines et
arguments identiques, \code{simuler\_base()} rend un jeu identique~; c'est
ce que vérifie, sur la plateforme où elle tourne, la batterie de
reproductibilité.

\textbf{Version du simulateur.} \code{VERSION\_SIMULATEUR} est consignée
au manifeste. Toute modification du module qui change la sortie de
\code{simuler\_base()} pour une entrée admise, fût-ce sur un seul jeu,
l'incrémente, et son commit est un changement de résultat sur les jeux,
avec son tableau avant~/ après\fsp; un refactorage sans effet garde la
version. La version garantit la sortie des jeux versionnés et de toute
entrée admise par la version courante. Les corrections de la revue finale
de la branche qui a introduit le simulateur font \textbf{exception
délibérée} à cette règle d'incrément~: les contrôles d'entrée refusent des
entrées autrefois admises, et la restriction des lignes éligibles à D01b
aux encaissements antérieurs à l'extraction change la sortie des entrées
encore admises dont un encaissement est daté de l'extraction (sa copie
n'est plus tirée). La décision du mainteneur du 9~novembre~2026 (M190)
admet ce changement sans incrément, aucun jeu n'ayant été publié hors de
cette branche, et les jeux versionnés n'étant pas changés (la batterie de
reproductibilité les rejoue à l'identique)~: la version reste~1. La sémantique des intensités de défauts de la loi vraie de
version~1 est figée~: ajouter un sous-cas demande une nouvelle version de
la loi vraie et du simulateur.

\textbf{Jeux versionnés.} Un jeu versionné est écrit par
\code{ecrire\_jeu()} sous \code{inst/extdata/jeux/<nom>/}, après le visa du
mainteneur, dont la date est portée au manifeste (décision M188), et n'est
jamais modifié en place. Un jeu d'une version antérieure du simulateur ou
du dictionnaire est retiré et remplacé par un jeu de numéro suivant
(décision M189). La batterie de reproductibilité rejoue chaque jeu
versionné à partir de ses seuls fichiers et le compare au jeu relu, champ à
champ, hors de la version du package qui l'a produit, contrôlée à part (au
plus la version installée)\fsp; elle compare aussi les empreintes, et
vérifie que les entrées du jeu sont celles du script de provenance
\code{tests/jeux\_v1.R}, qui porte en commentaire le statut et la
justification de chaque paramètre. Deux jeux sont versionnés~:
\code{v1-modeste-01}, jeu propre d'un portefeuille modeste à garanties
liées décès et arrêt de travail, et \code{v1-modeste-02}, même jeu propre
avec les douze défauts du catalogue, jeu de contrôle et non de calcul. Ils
servent aux tests fonctionnels, sans portée statistique.

\textbf{Affichage.} \code{format()} et \code{print()} décrivent un jeu par
des décomptes seuls, sans aucune somme.

Les constantes du module sont citées d'après le code
(\code{R/simulation.R})~:

\begin{center}
\small
\rowcolors{2}{tablerow}{white}
\begin{tabular}{@{}p{5.6cm}p{9.5cm}@{}}
\rowcolor{tableheader}
\textcolor{financeblue}{\textbf{Constante}} & \textcolor{financeblue}{\textbf{Valeur et objet}} \\
\code{FLUX\_SIMULATION} & \code{population}, \code{famille}, \code{sinistres}, \code{information}~: flux du jeu propre, dans leur ordre de tirage. \\
\code{GRAINES\_SIMULATION} & Les quatre flux, puis \code{defauts}~: noms de l'argument \code{graines}. \\
\code{VERSION\_LOI\_VRAIE} & Entier \code{1}~: version du format de la loi vraie. \\
\code{VERSION\_SIMULATEUR} & Entier \code{1}~: version du simulateur, consignée au manifeste. \\
\code{DEFAUTS\_SIMULES} & Catalogue des défauts injectables (rang, code, paramètre, table, lignes éligibles), douze lignes. \\
\code{COLONNES\_REGISTRE} & Colonnes du registre des défauts, dans leur ordre. \\
\code{SCHEMA\_VERITE} & Tables de vérité, leurs colonnes et leurs types (types du dictionnaire). \\
\end{tabular}
\end{center}


\partie{Partie XVII}{Module calculer\_provisions}

% =============================================================================
\section{Rôle et fonctions du module calculer\_provisions}
\label{sec:calculer-provisions}
% =============================================================================

Le module \code{calculer\_provisions} est l'orchestration~: il porte le
point d'entrée unique (section~\ref{sec:package}), la table des étapes, la
combinaison des postes à deux familles de calcul et le rejeu d'une fiche de
calcul. Il ne calcule aucun poste. Il peut appeler tous les modules, sauf
\code{lecture} et \code{simulation}. Ce module n'a pas encore de fonction.

\begin{center}
\small
\rowcolors{2}{tablerow}{white}
% concordance-fonctions : calculer_provisions
\begin{tabular}{@{}p{3.9cm}>{\raggedright\arraybackslash}p{3.0cm}p{8.2cm}@{}}
\rowcolor{tableheader}
\textcolor{financeblue}{\textbf{Fonction}} & \textcolor{financeblue}{\textbf{Arguments}} & \textcolor{financeblue}{\textbf{Rôle}} \\
\multicolumn{3}{@{}l@{}}{Aucune fonction.} \\
\end{tabular}
\end{center}

% =============================================================================
\section{Table des étapes}
\label{sec:etapes}
% =============================================================================

L'ordre du calcul est fixé par un tableau constant du module,
\code{ETAPES}, à quatre colonnes (\code{etape}, \code{module},
\code{fonction}, \code{condition}). L'orchestration le suit dans l'ordre,
par une séquence d'appels explicite, sans appel dynamique d'une fonction
désignée par son nom. Le tableau ci-dessous reproduit \code{ETAPES}\fsp; il
est vide tant que le module n'existe pas.

\begin{center}
\small
\rowcolors{2}{tablerow}{white}
\begin{tabular}{@{}p{2.6cm}p{3.6cm}p{4.2cm}p{4.2cm}@{}}
\rowcolor{tableheader}
\textcolor{financeblue}{\textbf{Étape}} & \textcolor{financeblue}{\textbf{Module}} & \textcolor{financeblue}{\textbf{Fonction}} & \textcolor{financeblue}{\textbf{Condition}} \\
\multicolumn{4}{@{}l@{}}{Aucune étape.} \\
\end{tabular}
\end{center}


\partie{Partie XVIII}{Tests}

% =============================================================================
\section{Batteries}
\label{sec:batteries}
% =============================================================================

Trois batteries se lancent depuis la racine du dépôt, dans cet ordre~:

\begin{encadre}
\small
\code{Rscript tests/test\_unitaires.R}\\
\code{Rscript tests/test\_reproductibilite.R}\\
\code{Rscript tests/concordance\_doc\_moteur.R -{}-strict}
\end{encadre}

\begin{itemize}[itemsep=2pt]
\item La \textbf{batterie unitaire} exécute d'abord les scripts de
  contrôle du dépôt (section~\ref{sec:controles}), chacun dans un processus
  \code{Rscript} distinct~: un script compte en échec si son statut de
  sortie n'est pas nul ou si sa sortie contient une ligne \code{[ECHEC]}.
  Elle installe ensuite le package dans une bibliothèque temporaire, le
  charge, puis exécute les fichiers de \code{tests/unitaires/}, un par
  module, plus ceux qui testent le harnais et la batterie de concordance.
\item La \textbf{batterie de reproductibilité} vérifie les aller-retours
  identiques du sérialiseur de listes, la graine locale et la restauration
  de l'état global du générateur, puis rejoue chaque jeu versionné de
  \code{inst/extdata/jeux/} à partir de ses seuls fichiers
  (section~\ref{sec:simulation-jeu}).
\item La \textbf{batterie de concordance} compare la présente documentation
  au code (section~\ref{sec:concordance}).
\end{itemize}

Les batteries unitaire et de reproductibilité testent le package installé, à
travers son espace de noms, jamais en chargeant les fichiers de \code{R/}
par \code{source()}\fsp; seules la batterie de concordance et les scripts
de contrôle lisent ces fichiers, sans les exécuter.

\textbf{Intégration continue.} Les trois batteries tournent, dans l'ordre
ci-dessus, dans le job «~Tests (R~4.3.3)~» de l'intégration continue du
dépôt, exécuté dans le conteneur \code{rocker/r-ver:4.3.3} épinglé par
digest~: c'est la \textbf{plateforme de référence}, sur
laquelle seront produites les références de non-régression. Avant les
batteries, le job imprime au journal \code{sessionInfo()},
\code{extSoftVersion()} et \code{La\_version()}. Le job «~Documents LaTeX~» compile la
charte graphique et le présent document par pdfLaTeX, en trois passes, et
échoue sur une erreur, une référence indéfinie ou une ligne trop longue
(\textit{Overfull})\fsp; il exécute ensuite le script de contrôle du
préambule. Aucun PDF n'est versionné.

% =============================================================================
\section{Scripts de contrôle du dépôt}
\label{sec:controles}
% =============================================================================

Six scripts de \code{tests/controles/} vérifient des règles du dépôt qui ne
sont pas des tests du package. Chacun s'exécute seul, depuis la racine du
dépôt (\code{Rscript tests/controles/<script>.R}), n'emploie que R de base,
imprime une ligne par constat (\code{[OK]}, \code{[ECHEC]}, et
\code{[AVERT]} pour le contrôle de Java), avec le fichier et la ligne en
cause, et sort avec un statut non nul au moindre échec. Un argument
facultatif désigne une autre racine, ce qui permet de démontrer qu'un
script refuse une copie du dépôt délibérément fautive (témoin négatif).

\begin{center}
\small
\rowcolors{2}{tablerow}{white}
\begin{tabular}{@{}p{5.3cm}p{9.8cm}@{}}
\rowcolor{tableheader}
\textcolor{financeblue}{\textbf{Script}} & \textcolor{financeblue}{\textbf{Ce qu'il vérifie}} \\
\code{liste\_blanche\_description.R} & Les champs \code{Depends}, \code{Imports}, \code{LinkingTo} et \code{Suggests} de \code{DESCRIPTION}, et les directives d'import de \code{NAMESPACE}, ne citent que des packages de priorité \code{base} (et R lui-même dans \code{Depends}). \\
\code{absence\_java.R} & Aucun package connu pour dépendre de Java n'est déclaré dans \code{DESCRIPTION}, ni présent dans la fermeture de ses dépendances parmi les packages installés\fsp; aucun fichier de \code{NAMESPACE}, \code{R/}, \code{app/}, \code{tests/} ou \code{.github/} ne charge un tel package ni n'installe un environnement Java. La liste des packages à Java est un jugement de mise en œuvre, écrit dans le script. \\
\code{graphe\_modules.R} & Les fichiers de \code{R/} sont parmi les quinze modules\fsp; chaque module n'appelle que les modules de sa colonne «~Peut appeler~» (section~\ref{sec:modules})\fsp; aucun appel dynamique (fonction désignée par son nom, accès à un espace de noms)\fsp; un nom de premier niveau n'est défini que dans un seul fichier. \\
\code{frontiere\_calcul\_affichage.R} & Aucun élément de Shiny dans \code{R/}\fsp; aucun réglage du générateur pseudo-aléatoire (\code{set.seed()}, \code{RNGkind()}, \code{RNGversion()}, \code{.Random.seed}) hors de \code{R/outils.R}\fsp; aucun appel de \code{warning()}, \code{message()}, \code{cat()}, \code{print()}, \code{writeLines()}, \code{writeChar()}, \code{packageStartupMessage()}, \code{Sys.time()}, \code{Sys.Date()} ni \code{date()}, hors du corps des fonctions de premier niveau nommées \code{print.*} et \code{format.*}\fsp; code ASCII hors commentaires\fsp; si \code{app/} existe, aucun appel aux fonctions internes du package ni à une liste de fonctions de calcul. \\
\code{concordance\_dictionnaire.R} & Le dictionnaire lisible par machine (section~\ref{sec:dictionnaire-csv}) concorde avec le texte visé \code{docs/specifications/dictionnaire-base-v1.md}, qui fait foi~: le dossier contient exactement les sept CSV et \code{dictionnaire.txt}\fsp; la version de \code{dictionnaire.txt} est celle du texte et \code{VERSION\_DICTIONNAIRE\_ATTENDUE}, sa date de visa celle de l'historique du texte, ses empreintes celles des CSV\fsp; chaque CSV est égal, ligne à ligne, champ à champ et dans le même ordre, au tableau correspondant du texte\fsp; les codes de contrôle sont disjoints des codes de défaut, et ceux qui sont des motifs de refus d'appel sont exactement \code{colonne\_absente}, \code{colonne\_inconnue} et \code{type\_incompatible}. Les limites de détection sont écrites en tête du script. \\
\code{preambule\_latex.R} & Le préambule de chaque document LaTeX de \code{docs/latex/} est identique à celui de la charte graphique, hors des blocs bornés par «~À ADAPTER~», dont le nombre, l'ordre et l'objet doivent être ceux de la charte. \\
\end{tabular}
\end{center}

Ces scripts sont des \textbf{contrôles statiques}~: ils analysent le texte
des fichiers, sans exécuter le package. Ils ne prouvent pas l'absence de
toute violation~: les cas résiduels de faux négatifs et de faux positifs
relevés par l'audit de leur première version sont suivis par l'issue~27 du
dépôt.

% =============================================================================
\section{Harnais}
\label{sec:harnais}
% =============================================================================

Le harnais (\code{tests/harnais.R}) est l'outillage de test du projet, sans
le package \code{testthat}. Chaque batterie le source\fsp; il exige que les
avertissements soient convertis en erreurs (\code{options(warn = 2)}), faute
de quoi il arrête la batterie. Il fournit~:

\begin{itemize}[itemsep=2pt]
\item \code{section()}, qui ouvre une section de la sortie\fsp;
\item \code{verifier()}, qui vérifie une condition\fsp;
\item \code{verifier\_egal()}, qui vérifie une égalité, exacte ou à une
  tolérance relative déclarée\fsp;
\item \code{verifier\_refus()}, qui vérifie qu'un refus d'appel est levé
  avec le motif attendu\fsp;
\item \code{attendu\_echec()}, pour un défaut connu, avec renvoi à son
  issue\fsp;
\item \code{bilan()}, qui imprime le décompte et termine en échec s'il y a
  un échec ou un succès inattendu\fsp;
\item \code{installer\_provie()}, qui installe et charge le package\fsp; si
  l'installation échoue, il imprime la sortie de \code{R CMD INSTALL}, puis
  une ligne \code{[ECHEC]}, et arrête la batterie avec le statut~1.
\end{itemize}

Chaque vérification imprime une ligne~: \code{[OK]}, \code{[ECHEC]},
\code{[ATTENDU]} (échec attendu, avec l'issue) ou \code{[SUCCES INATTENDU]}.
Un échec attendu qui réussit fait échouer la batterie~: la marque est alors
retirée pour en faire un test ordinaire.

% =============================================================================
\section{Concordance entre la documentation et le code}
\label{sec:concordance}
% =============================================================================

La batterie \code{tests/concordance\_doc\_moteur.R} lit les tableaux de
fonctions du présent document et les compare aux fonctions de premier
niveau de chaque fichier \code{R/<module>.R}, hors auxiliaires dont le nom
commence par un point. Elle relève comme écarts~:

\begin{itemize}[itemsep=2pt]
\item une fonction du code absente du tableau de son module, ou une
  fonction citée absente du code\fsp;
\item une fonction citée deux fois, deux tableaux pour un même module, un
  module inconnu\fsp;
\item un module qui a des fonctions mais aucun tableau\fsp;
\item un fichier de \code{R/} hors des quinze modules\fsp;
\item une fonction exportée par le package qui n'est citée dans aucun
  tableau.
\end{itemize}

\textbf{Valeurs des constantes.} La batterie lit aussi les tableaux de
constantes, reconnus à la première cellule de leur en-tête,
«~Constante~»\fsp; chacun appartient au module dont le tableau de
fonctions le précède. Elle évalue les constantes de premier niveau de
\code{R/<module>.R} dont le nom est en majuscules, sans charger le
package, et compare à leur valeur quatre formes littérales écrites dans
la deuxième ou la troisième cellule d'une rangée~:

\begin{itemize}[itemsep=2pt]
\item une chaîne, \code{"texte"} en tête de la deuxième cellule\fsp;
\item un entier, «~entier~\code{N}~» en tête de la deuxième cellule, pour
  une constante de type entier\fsp;
\item un décompte, «~quatorze lignes~», «~dix-neuf motifs~», en chiffres
  ou en lettres jusqu'à vingt, comparé à la longueur de la constante ou à
  son nombre de lignes\fsp;
\item une liste de valeurs écrites sans guillemets, comparée dans
  l'ordre.
\end{itemize}

Une valeur rédigée autrement n'est pas lue~: la batterie nomme la
constante «~sans valeur littérale lue~», sans écart. Elle relève comme
écarts une valeur différente de celle du code, une constante citée absente
de son module, non évaluable ou citée deux fois, un tableau de constantes
qui ne suit aucun tableau de fonctions, un décompte composé en lettres
(qu'elle ne sait pas lire), une rangée mal fermée ou plus large que
l'en-tête, et toute chaîne entre guillemets qu'aucune forme n'a lue
(décision du mainteneur du 9~novembre~2026, M207)~: une version citée au
fil d'une phrase, une chaîne en troisième cellule, une liste écrite avec
guillemets. Une chaîne entre guillemets écrite dans une cellule de valeur
est donc toujours comparée au code ou signalée. Les constantes internes,
dont le nom commence par un point, ne sont pas lues\fsp; elles se citent
dans le texte, d'après le code.

Avec l'option \code{-{}-strict}, tout écart, comme l'absence du document,
fait échouer la batterie\fsp; sans elle, les écarts sont des avertissements.
La convention d'écriture des tableaux qu'elle lit est fixée par
\code{docs/latex/CONVENTIONS.md}.

\end{document}
